% GT class project paper for JFM

\documentclass{jfm}

\usepackage{mathptmx}
\usepackage{natbib}
\usepackage{graphicx}
\usepackage{amsmath}
\usepackage{amssymb}
\usepackage{verbatim}
%\usepackage{hyperref}

\voffset .4in
\hoffset .3in

\renewcommand{\baselinestretch}{1.4}          % enlarged line skip

\newcommand{\ie}  {\textit{i.e., }}
\newcommand{\eg}  {\textit{e.g., }}
\newcommand{\ca}  [1]{\mathcal{#1}}
\newcommand{\kb}  {\boldsymbol{k}}
\newcommand{\bs}  [1]{\boldsymbol{#1}}
\newcommand{\rms} {\mathrm{rms}}
\newcommand{\del} {\nabla}
\newcommand{\p}   {\partial}
\newcommand{\pp}  [2]{\frac{\p #1}{\p #2}}
\newcommand{\dd}  [2]{\frac{d #1}{d #2}}

\title[Turbulent Diffusion in the Geostrophic Inverse Cascade]
    {Turbulent Diffusion in the Geostrophic Inverse Cascade}

\author[Smith, Boccaletti, Henning, Marinov, Tam, Held and Vallis]
    {
    K.\ns S.\ns S\ls M\ls I\ls T\ls H,\ns
    G.\ns B\ls O\ls C\ls C\ls A\ls L\ls E\ls T\ls T\ls I,\ns \\
    C.\ns C.\ns H\ls E\ls N\ls N\ls I\ls N\ls G,\ns
    I.\ns M\ls A\ls R\ls I\ls N\ls O\ls V,\ns
    C.\ns Y.\ns T\ls A\ls M,\ns \\
    I.\ns M.\ns H\ls E\ls L\ls D,\ns
    G.\ns K.\ns V\ls A\ls L\ls L\ls I\ls S
    }

\affiliation{GFDL/Princeton University, Princeton, New Jersey}

\date{\today}

\begin{document}

\maketitle

\begin{abstract}
  Motivated in part by the problem of large-scale lateral
  turbulent heat transport in the earth's atmosphere and oceans,
  and in part by the problem of turbulent transport itself,
  we seek to better understand the transport of a
  passive tracer advected by various types of fully developed
  two-dimensional turbulence. The types of turbulence considered
  correspond to various relationships between the streamfunction
  and the advected field. Each type of turbulence considered
  possesses two quadratic invariants and each can develop an
  inverse cascade. These cascades can be modified or halted, for
  example, by friction, a background vorticity gradient or a mean
  temperature gradient. We focus on three physically-realizable
  cases: classical two-dimensional turbulence, surface
  quasi-geo\-strophic turbulence, and shallow water
  quasi-geo\-strophic turbulence at scales large compared to the
  radius of deformation. In each model we assume that tracer
  variance is maintained by a large-scale mean tracer gradient
  while turbulent energy is produced at small scales via random
  forcing, and dissipated by linear drag.  We predict the spectral
  shapes, eddy scales and equilibrated energies resulting from the
  inverse cascades, and use the expected velocity and length
  scales to predict integrated tracer fluxes.

  When linear drag halts the cascade, the resulting diffusivities
  are decreasing functions of the drag coefficient, but with
  different dependencies for each case.  When $\beta$ is
  significant, we find a clear distinction between the tracer
  mixing scale, which depends on $\beta$ but is nearly
  independent of drag, and the energy containing (or jet) scale,
  set by a combination of the drag coefficient and $\beta$.  Our
  predictions are tested via high resolution spectral
  simulations.  We find in all cases that the
  passive scalar is diffused down-gradient with a diffusion
  coefficient that is well predicted from estimates of mixing
  length and velocity scale obtained from turbulence
  phenomenology.
\end{abstract}

\section{Introduction}
\label{intro:sec}

This paper is concerned with the problem of transport of a
passive tracer in a turbulent flow. In addition to its intrinsic
interest, the problem is relevant to the issue of meridional heat
transport in the earth's atmosphere and oceans. In mid-latitudes
much of the atmospheric transport is effected by large-scale
eddies that are well described by the quasi-geostrophic equations
of motion \citep[e.g.,][]{Pedlosky87}. In the classical two-layer
model of geostrophic turbulence
\citep[e.g.,][]{Rhines77,Salmon80}, at scales larger than the
deformation radius the baroclinic streamfunction is nearly
passively advected by the barotropic streamfunction, resulting in
a down-scale cascade of its variance. But energy in the
barotropic mode itself cascades to large scales. Since temperature is
proportional to the baroclinic streamfunction, its nearly-passive
advection by the barotropic flow determines the heat transport,
and this down-gradient flux of heat in turn determines the
extraction of energy from the mean flow \citep{Larichev95}.  To
understand the heat transport properties of this system we must
at least, then, understand the transport properties of a passive
tracer being advected by a turbulent flow in which energy is
cascading to larger scales, and that is the subject of this paper.

A familiar setting in which to explore this problem is that of
conventional two-di\-men\-sion\-al turbulence [see e.g.,
\citet{Kraichnan80}, \citet{Vallis93b} and \citet{danilov00} for
reviews]. In such flows, vorticity $\zeta$ is advected by a flow
defined by a streamfunction $\psi$ linked to the vorticity, in Fourier
transform space, by $\zeta = -|\kb|^2 \psi$.  We will find it useful
to consider more general relationships between streamfunction and
advected field of the form $\xi = -|\kb|^{\alpha} \psi$. This is not
only a formal device, for a number of physical systems can be realized
by different choices of $\alpha$. For example, in surface
quasi-geostrophic dynamics \citep[\eg][]{Held95a}, streamfunction
$\psi$ and advected quantity $\xi$ are related by $\xi = -|\kb| \psi$,
and the equations represent the advection of temperature along a
surface bounding a constant potential vorticity interior, so providing
(for example) a simple model for edge waves on the tropopause or
temperature advection near the Earth's surface. Large-scale
quasi-geostrophic flow (quasi-geostrophic dynamics at scales large
compared to the radius of deformation), may be a particularly useful
model for Jovian dynamics, presuming that small-scale convection
constitutes a source of turbulent energy.  The latter case corresponds
to, as we will see, yet a third relationship between advected scalar
and associated streamfunction.

The dynamics of such generalised two-di\-men\-sion\-al turbulence was
explored by \citet{Pierrehumbert94} with an emphasis on the forward
cascade properties of the flow. Our focus is on the inverse cascade
and the transport of a passive tracer by such flows. Thus, we consider
fluids that are stirred at small scales by a random forcing,
generating an inverse cascade, and which advect a passive tracer whose
variance is maintained by a large-scale gradient. Most of the
transport of the tracer thus occurs at scales within the inverse
cascade of the turbulent fluid, just as for heat transport in
baroclinic turbulence. We derive general expressions for the spectral
shape of the inverse cascade in the presence of friction, as well as
estimates of the halting scale, equilibrated energy level and tracer
diffusivity. We then test these expressions numerically for the three
cases of geostrophic turbulence mentioned above --- conventional
two-dimensional vorticity dynamics, surface quasi-geostrophic
dynamics, and large-scale quasi-geostrophic dynamics --- each of which
have a different relationship between streamfunction and advected
scalar. In conventional two-dimensional turbulence we also explore the
effects of differential rotation, and we will see that the stopping
scale of the inverse cascade is not necessarily the same as the mixing
length of a passive scalar; we offer, and test, expressions for these
two scales.

The theory presented in the paper is largely that of classical
homogeneous turbulence, and accordingly we rely upon
Kolmogorov-Kraichnan phenomenology, the presumption of well-defined
eddy scales and magnitudes, and implicitly upon the convergence of
flows to statistically steady states.  Kolmogorov-Kraichnan
phenomenology has been investigated for two-dimensional flows by many
authors, and is generally found to hold best for flows that are
largely free of coherent structures and devoid of intermittency.
Examples of this can be found in, \eg \citet{Maltrud91} and
\citet{Sukoriansky99} for the inverse cascade and \citet{Oetzel97} and
(at much higher resolution) \citet{Lindborg00} for the direct cascade.
In the simulations we present, the flows are forced by small-scale
random noise, which appears to minimize the production of coherent
vortices in the inverse cascade of standard two-dimensional
turbulence.  Moreover the inverse cascade is, in all cases, halted
\emph{before} the up-scale cascading invariant reaches the domain
scale, preventing vortex condensation \citep{SmithYakhot93}.
Nevertheless, in some cases investigated, coherent structures do form,
and we will point out where these deviations from phenomenology occur.

Such phenomenology is sometimes questioned in light of simulations of
two-dimensional turbulence forced at small scales reported by
\citet{borue94}.  In those simulations it is found that while after a
few eddy turn-around times the energy spectrum possesses a -5/3 slope,
at much later times the spectrum evolves to a -3 slope, accompanied by
coherent vortices in the flow.  However, Borue employed an inverse
hyper-viscosity (or `hypo'-viscosity) to absorb energy at large
scales, rather than the more traditional and geophysically relevant
linear drag.  The former method of energy dissipation has recently
been shown \citep{danilov01a} to cause the steepening of the energy
spectrum in the inverse cascade observed by Borue, while linear drag
merely tapers the -5/3 spectrum at large scales, so long as the drag
is large enough to prevent significant energy from reaching the domain
scale.

The paper is organized as follows. In \S \mbox{} \ref{gen2d:sec} we introduce
the basic formalism of generalised two-dimensional turbulence, and in
\S \ref{spectra:sec} we discuss the fundamental spectral properties of
such flows and calculate their corresponding inertial range spectra
using Kolmogorov-Kraichnan phenomenology.  The dynamics of a passive
tracer in an inverse cascade of generalised two-dimensional turbulence
are discussed in \S \ref{tracer:sec}.  In \S \ref{halting:sec} we
heuristically derive a shape for the spectrum of the inverse cascade
in the presence of drag, which in turn yields an estimate of the
halting scale of the cascade. In \S \ref{beta:sec} we add to the
picture a mean vorticity gradient (the $\beta$-effect). This produces
zonal jets, and we suggest expressions for both the jet scale and the
eddy mixing scale. Finally, in \S \ref{numerics:sec}, we report on
numerical tests of our scaling predictions, and conclude in \S
\ref{disc:sec}.  An argument for cascade directions is given in
appendix A and some details of the numerical model are relegated to
appendix B.

\section{Two-dimensional flow: a generalised formalism}
\label{gen2d:sec}

Consider the two-dimensional advection equation for a conserved scalar
field $\xi$
\begin{equation}
\pp{\xi}{t} + J(\psi,\xi) = 0,
\label{gen2d:eq}
\end{equation}
where $\psi$ is the advecting streamfunction and $J(A,B) = \p_x A\p_y
B - \p_y A\p_x B$ is the Jacobian operator. The fluid velocity is
given by $\boldsymbol{u} = \hat{k}\times\nabla\psi$.

In ordinary two-dimensional turbulence $\xi = \del^2 \psi$.
Restricting attention to a doubly-periodic domain, in Fourier
transform space, this is equivalent to the relationship between
Fourier components
\begin{equation}
\xi_{\kb} = -k^2\psi_{\kb},
\label{coupling1:eq}
\end{equation}
where $k=|\kb|$ is the isotropic two-dimensional wavenumber. More
generally we can consider the relationship between $\xi_{\kb}$ and
$\psi_{\kb}$ discussed by \citet{Pierrehumbert94},
\begin{equation}
\xi_{\kb} = -k^{\alpha}\psi_{\kb},
\label{coupling:eq}
\end{equation}
where $\xi$ shall be referred to as the generalised vorticity
field. The independent variables $x,y$ have units of length $[L]$, and
$t$ has units of time $[T]$.  Therefore, in order for \eqref{gen2d:eq}
to be dimensionally balanced, $\psi \sim [L^2T^{-1}]$ and $\xi \sim
[L^{2-\alpha}T^{-1}]$.  Three particular values of $\alpha$ lead to
geophysically relevant equations: $\alpha=2$, $\alpha = 1$ and $\alpha
= -2$.

\bigskip
\noindent\textit{$\alpha = 2$: Two-dimensional vorticity dynamics}

\begin{narrower}

\noindent If $\alpha=2$ we obtain the familiar two-dimensional Euler, or
ba\-ro\-tro\-pic vor\-ticity equation, with $\xi = \zeta$ and $\zeta \sim
[T^{-1}]$. In this paper we shall refer to these dynamics as
two-dimensional vorticity (TDV) dynamics.

\end{narrower}

\bigskip
\noindent\textit{$\alpha = 1$: Surface quasi-geostrophic dynamics}

\begin{narrower}

\noindent When $\alpha=1$ we obtain surface quasi-geostrophic (SQG)
dynamics \citep[\eg][and references therein]{Held95a}. These equations
describe the evolution of surface temperature perturbations bounding a
constant potential vorticity interior, within the quasi-geostrophic
equations.

\end{narrower}

\bigskip
\noindent\textit{$\alpha = -2$: Large-scale quasi-geostrophic dynamics}

\begin{narrower}

\noindent The case $\alpha=-2$ corresponds to a rescaled shallow water
quasi-geostrophic equation in the asymptotic limit of length scales
large compared to the deformation scale [this system was
investigated by \citet{Larichev91}]. We term this type of flow
\textit{large-scale quasi-geostrophic (LQG) dynamics.} Consider the
shallow water qua\-si\-geo\-stro\-phic equations
\begin{equation}
\pp{q}{t} + J(\psi,q) = 0, ~~~ q = (\del^2 - \lambda^2) \psi
\label{qg:eq}
\end{equation}
where $\lambda$ is the deformation wavenumber.  In Fourier space the
expression for the potential vorticity $q$ is
\begin{equation}
q_{\kb} = -(k^2 + \lambda^2) \psi_{\kb}
\label{qk:eq}
\end{equation}
For wavenumbers such that $k\gg \lambda$, we recover TDV, but when
$k\ll \lambda$, \eqref{qg:eq} becomes
\begin{equation}
-\lambda^2\pp{\psi}{t} + J(\psi,\del^2 \psi) = 0.
\label{lqg1:eq}
\end{equation}
Because $\lambda$ only appears in combination with the time derivative,
we can rescale time
\begin{equation}
\tau = t\lambda^{-2},
\label{tau:eq}
\end{equation}
and make the substitution of variables
\begin{equation}
\xi =  \psi ~~~ \textrm{and}~~~ \Psi = \del^2 \psi
\label{lqgvar:eq}
\end{equation}
so that \eqref{lqg1:eq} becomes
\begin{equation}
\pp{\xi}{\tau} + J(\Psi,\xi) = 0,~~~\Psi = \del^2 \xi.
\label{lqg3:eq}
\end{equation}
In Fourier space the relation between $\xi$ and $\Psi$ is
\begin{equation}
\xi_{\kb} = k^{-2}\Psi_{\kb}
\label{xik:eq}
\end{equation}
which has the same form as \eqref{coupling:eq} with $\alpha=-2$.

\end{narrower}

\bigskip

Nonlinear advection is present in all of these geophysical cases, thus
turbulence and turbulent cascades can develop.  The character of the
turbulence --- \eg the spectral shapes, eddy scales, cascade
directions, statistical amplitudes, coherent structures and rates of
turbulent development --- will be distinct in each case.  Some aspects
of these distinctions were discussed by \citet{Pierrehumbert94}, who
argued, for example, that the enstrophy cascade should be dominated by
local strain for $\alpha<2$, and that local scaling predictions should
break down for $\alpha>2$.  Furthermore, they note, the inverse
cascade --- our primary concern in this paper --- should remain local
for all $\alpha<4$; the three cases discussed above fall into this
category.

The fact that a single generalised formalism describes three cases of
geophysical interest by changing the value of a single parameter
provides strong motivation in itself to investigate the properties of
the general system. We may also regard the generalised turbulence
formalism, however, as a convenient shorthand with which we can derive
common aspects of the flows without repetition.

In order to make quantitative statements about scales and
energies (a necessary prerequisite to predicting the statistics
of passive tracer advection) we first consider the Kolmogorov
spectra for the generalised dynamics.  Some of what follows in
the next section was derived by \citet{Pierrehumbert94}, but we
include it for clarity, notation and as a reference basis for the
rest of the paper.

\section{Spectral properties of generalised two-dimensional flow}
\label{spectra:sec}

Because $\xi$ is conserved on parcels, the integral over a periodic
domain (or an enclosed domain with no-normal-flow boundary conditions)
of any function of $\xi$ is conserved. Thus, there are an infinite
number of invariants. However, just as for TDV flow, two quadratic
invariants determine the cascade directions. From \eqref{gen2d:eq}
these are
\begin{equation}
E_G\equiv-\frac{1}{2}\overline{\psi\xi},~~~
Z_G\equiv\frac{1}{2}\overline{\xi^2},
\label{invar:eq}
\end{equation}
where the overline indicates a horizontal average.  We will refer to
$E_G$ as the generalised energy, and to $Z_G$ as the generalised
enstrophy. The invariants have units $E_G \sim [L^{4-\alpha}T^{-2}]$
and $Z_G \sim [L^{4-2\alpha}T^{-2}]$.  [Note that for LQG, $T$ is the
rescaled time from \eqref{tau:eq}]. The spectra of the two quadratic
invariants are connected by a simple relationship involving scale
only. To see this, we define the isotropic spectra $\ca{A}(k)$ of the
positive definite quantity $A$ such that
\begin{equation}
A = \int_0^{\infty}\ca{A}(k) \, dk,
\label{Adef:eq}
\end{equation}
where $k$ is the isotropic wavenumber, and $\ca{A}(k)$ is either
$\ca{E}_G(k)$ or $\ca{Z}_G(k)$ (thus the dimensions of $\ca{E}_G(k)
\sim [L][E_G]$ and $\ca{Z}_G(k) \sim [L][Z_G]$). The coupling
relationship \eqref{coupling:eq}, together with Parseval's theorem,
implies that
\begin{equation}
\ca{Z}_G(k) = k^{\alpha}\ca{E}_G(k).
\label{EZ:eq}
\end{equation}

Well-known arguments \citep{Batchelor53} lead one to expect that,
in TDV flow, energy will cascade to larger scale and enstrophy
will cascade to smaller scale. These arguments are based on the
presumption that the distribution of the spectra of each
invariant spreads in wavenumber space.  In appendix A we extend
such arguments to the case with arbitrary $\alpha$ and show that,
for $\alpha>0$, generalised energy cascades to larger scale and
generalised enstrophy cascades to smaller scales.  For
$\alpha<0$, however, it turns out that generalised energy
cascades to \textit{small} scales, and generalised enstrophy
cascades to \textit{large} scale.

If there is a spectrally localised source of variance of $E_G$ and
$Z_G$, then in all the above cases we expect there will be a
transfer of generalised energy and generalised enstrophy in
opposite directions, away from the source region, just as in TDV
dynamics. If frictional effects only act at some spectral
distance from the source then we might expect inertial ranges to
form between source and sink.  If this is the case one can derive
the spectral slopes of these putative inertial ranges.

Between source and sink the spectral flux $\epsilon_E$ or $\epsilon_Z$
associated with the generalised energy or enstrophy will be constant,
producing an inertial range.  We assume that over their respective
inertial ranges, the local spectral density $\ca{E}_G(k)$ and
$\ca{Z}_G(k)$ are functions only of the local spectral flux of $E_G$
or $Z_G$ and of the scale itself --- the locality hypothesis of
\citet{Kolmogorov41}. Thus
\begin{equation}
\epsilon_A = \frac{k\ca{A}(k)}{T_A(k)} = \textrm{constant},
\label{epsA:eq}
\end{equation}
where $A$ is $E_G$ or $Z_G$ and $T_A(k)$ is a local timescale.  From
dimensional considerations
\begin{subequations}
\begin{align}
T_E(k) = [k^{5-\alpha}\ca{E}_G(k)]^{-1/2}, \label{TE:eq} \\
T_Z(k) = [k^{5-2\alpha}\ca{Z}_G(k)]^{-1/2}. \label{TZ:eq}
\end{align}
\end{subequations}
Using the expression for the flux \eqref{epsA:eq} we obtain
\begin{subequations}
\begin{align}
\ca{E}_G(k) = \ca{C}_E\epsilon_E^{2/3}k^{(\alpha-7)/3},
\label{Espec:eq} \\
\ca{Z}_G(k) = \ca{C}_Z\epsilon_Z^{2/3}k^{(2\alpha-7)/3},
\label{Zspec:eq}
\end{align}
\end{subequations}
where $\ca{C}_E$ and $\ca{C}_Z$ are the Kolmogorov constants for the
respective cascades of $E_G$ and $Z_G$.  Because the dynamics for each
type of turbulence are different, $\ca{C}_E$ and $\ca{C}_Z$ may be
expected to depend on $\alpha$.  Note also that one can find the
spectrum $\ca{Z}_G(k)$ in the $E_G$-range or $\ca{E}_G(k)$ in the
$Z_G$-range via \eqref{EZ:eq}.

As a useful shorthand, we combine these results by writing the
spectrum of a generalised invariant as
\begin{equation}
\ca{A}(k) = \ca{C}\epsilon^{2/3}k^{-\gamma},
\label{Aspec:eq}
\end{equation}
where the exponent $\gamma$ is
\begin{equation}
\gamma = \left\{
    \begin{array}{ll}
       (7-\alpha)/{3},   &  A = E_G \\
      (7-2\alpha)/{3},   &  A = Z_G.
    \end{array}
\right.
\label{gam:eq}
\end{equation}
Note that as defined, $\gamma\geq 1$ for all cases considered in this
paper.  For future reference, using \eqref{epsA:eq} and
\eqref{Aspec:eq}, the generalised eddy timescale is then
\begin{equation}
T_A(k) = [k^{3\gamma-2}\ca{A}(k)]^{-1/2}
= \ca{C}^{-1/2}\epsilon^{-1/3}k^{1-\gamma}.
\label{TA:eq}
\end{equation}

Schema of the expected cascade directions and spectral slopes for each
of the three cases considered in this paper can be found in figure
\ref{schema:fig}.  Also noted are the dimensional estimates of the
linear drag-induced stopping scale (discussed in \S
\ref{halting:sec}), and the spectra of the velocity variance $\ca{E}$
for each case.  The latter is necessary since rms velocities shall be
needed to predict diffusivities in \S \ref{tracer:sec}.  In order to
calculate the spectral slope of the velocity variance one must
multiply $\ca{A}(k)$, the upscale cascading invariant, by the power of
$k$ necessary to give a spectrum with dimensions $[L^3T^{-2}]$, a
power which varies from case to case since $\ca{A}(k)$ has dimensions
which depend on $\alpha$.  One must be careful in the LQG case, since
in this case time has been re-dimensioned via \eqref{tau:eq}.
Particulars of the velocity estimates will be addressed for each case
separately as necessary in the context of the numerical calculations
discussed in \S \ref{numerics:sec}.

\section{Passive tracer dynamics}
\label{tracer:sec}

Suppose now that the streamfunction $\psi$ determined by the advection
equation \eqref{gen2d:eq} and the coupling \eqref{coupling:eq} also
advects a passive tracer $\phi$, for which the equation is
\begin{equation}
\pp{\phi}{t} + J(\psi,\phi) = \ca{D}_{\phi},
\label{tracer:eq}
\end{equation}
where $\ca{D}_{\phi}$ represents dissipation, presumed to occur at
small scales.  [In the LQG case, the streamfunction advecting the
tracer is the QG streamfunction $\psi$, which determines the
velocities, rather than the redefined streamfunction of
\eqref{lqgvar:eq}.] The tracer variance
\begin{equation}
P \equiv \frac{1}{2}\overline{\phi^2}
\label{P:eq}
\end{equation}
is a conserved quantity of the tracer field in the absence of
dissipation.  In the inertial range of $A$ --- some conserved quantity
of the flow field --- the spectral flux $\chi$ of tracer variance is
\begin{equation}
\chi = \frac{k\ca{P}(k)}{T_E(k)} = \textrm{constant}
\label{epsP:eq}
\end{equation}
where $\ca{P}(k)$ is the spectrum of the tracer variance and
$T_E(k)$ is the timescale associated with the kinetic energy in
the inertial range of $A$.  In general there is only one
spectrally local timescale available, namely \eqref{TA:eq}.  In
this case the tracer variance spectrum is
\begin{equation}
\ca{P}(k) = \ca{K}\chi\epsilon^{-1/3}k^{-\gamma},
\label{PE:eq}
\end{equation}
where $\ca{K}$ is the Kolmogorov constant for the tracer. In
general, then, the tracer variance follows the same power law as
the spectrum of the advecting flow invariant, in that invariant's
inertial range.  For example, in the case $\alpha=2$ (TDV), tracer
variance has the slope $-5/3$ in the energy inertial range, and
slope $-1$ in the enstrophy inertial range.

When LQG describes the dynamics, however, these predictions should
only hold if the tracer is advected by the redefined streamfunction.
This is not the relevant physical situation.  The variance of the
tracer as given by \eqref{tracer:eq} requires that we specify the
timescale $T_E(k)$ in terms of physical scales, which re-introduces
the deformation scale into the calculation.  The up-scale cascading
invariant is then associated with the available potential energy (APE)
and the direct cascading invariant is the kinetic energy (KE).  The
latter has a spectral slope of $k^{-5/3}$ in the inertial range of the
APE, yielding an expected tracer variance spectral slope of $k^{-5/3}$
as well.

Suppose now that tracer variance is generated by the presence of a
fixed large-scale meridional gradient and absorbed by small-scale
dissipation. Specifically, we consider
\begin{equation}
\pp{\phi}{t} + J(\psi,\phi) + \overline{\theta} v = \ca{D}_{\phi}.
\label{tracermg:eq}
\end{equation}
where $v=\p\psi/\p x$ is the meridional velocity of the advecting
flow and $\overline{\theta} \equiv \p\bar{\phi}/\p y$ is the mean
tracer gradient.  Since \eqref{tracermg:eq} is linear in $\phi$,
we can rescale the equation such that $\overline{\theta}=1$.
Multiplying \eqref{tracermg:eq} by $-\phi$ and averaging over
space yields
\begin{equation}
\dd{P}{t} = \overline{\phi v} - \overline{\phi \ca{D}_{\phi}},
\label{Pbal:eq}
\end{equation}
where $P$ is the tracer variance of \eqref{P:eq}.  The meridional flux
of tracer is the explicit source of variance in the tracer field, and
can be written
\begin{equation}
\overline{\phi v} = -D,
\label{td:eq}
\end{equation}
where $D$ is the diffusivity, which has dimensions $[L^2T^{-1}] =
[VL]$.  If the advecting meridional velocity field is sharply peaked
at some scale $k_0$ and has a rms value $V_{\rms}$, then a standard
down-gradient mixing hypothesis \citep[\eg][chapter 8]{Tennekes94}
yields an expected diffusivity
\begin{equation}
D \simeq V_{\rms}k_0^{-1}.
\label{diff:eq}
\end{equation}
One can make a crude estimate of the tracer flux spectrum by
decomposing the left hand side of \eqref{td:eq} and estimating the two
fields separately, yielding
\begin{equation}
\ca{D}(k) \sim \ca{E}^{1/2}(k)\ca{P}^{1/2}(k),
\label{Dspec:eq}
\end{equation}

We shall derive specific diffusivity estimates for each of our
cases and test their accuracy numerically.  In order to do
so, however, we must first derive estimates for $k_0$ and
$V_{\rms}$.

\section{Cascade reduction by linear drag}
\label{halting:sec}

An inverse cascade will be diminished or halted when the
upscale-cas\-ca\-ding in\-va\-ri\-ant either reaches the domain scale,
or the scale at which some competing process such as friction
dominates non-linear advection.  If some mechanism halts the cascade
of whichever invariant is cascading upscale at wavenumber $k_0$, and
if the eddies are forced at wavenumber $k_f$, then we must presume
that
\begin{equation}
k_f>k_0.
\label{kfkl:eq}
\end{equation}
If the spread between these scales is large, an inertial range may
exist between them.

As the eddies cascading through the system become larger, they also become
slower.  In particular, we presume an eddy timescale given by
\eqref{TA:eq}.  If the mechanism competing against nonlinear advection has
a timescale ($T_{\mathrm{halt}}$) that is independent of scale, or
increases less rapidly than $T_{\mathrm{eddy}}$ with increasing scale, then
at some point in the cascade, $T_{\mathrm{eddy}} \sim T_{\mathrm{halt}}$.
This expression can presumably be used to find the scale $k_0$ at which the
cascade is halted, analogous to the derivation of the Kolmogorov
micro-scale at which molecular viscosity overcomes the direct cascade.

Consider the case of a linear vorticity drag with the equation of motion
\begin{equation}
\pp{\xi}{t} + J(\psi,\xi) = F - r\xi,
\label{gen2dr:eq}
\end{equation}
where $F$ is the forcing. Physically, the term $r \xi$ represents a
drag due to Ekman friction in TDV, or thermal damping in SQG or
LQG. Note that $r \sim [T^{-1}]$: the drag coefficient has dimensions
of an inverse time for all couplings $\alpha$.  For TDV with linear
drag, \citet[p. 122]{Arbic00} proves that \eqref{kfkl:eq} must hold
for all values of drag coefficient.

If the time-averaged generalised energy (enstrophy) injection
rate $g$ is known and $\alpha>0$ ($\alpha<0$), the equilibrated
generalised energy (enstrophy) level can be calculated exactly.
To see this, multiply \eqref{gen2dr:eq} by $-\psi$ ($-\xi$) and
integrate over space to form the generalised energy (enstrophy)
budget equation
\begin{equation}
\dot{A} = g - 2rA,
\label{Atot:eq}
\end{equation}
where $g$ is the generation rate
\begin{equation}
g = \left\{
    \begin{array}{ll}
       -\overline{\psi F},   &  \alpha>0 \\
       -\overline{\xi F},    &  \alpha<0.
    \end{array}
\right.
\label{gdef:eq}
\end{equation}
In steady state
\begin{equation}
A = A_r = \frac{g}{2r},
\label{Ar2:eq}
\end{equation}
thus we know in advance the total generalised energy (enstrophy) of
the statistically steady flow for $\alpha>0$ ($\alpha<0$).  We have
neglected the effects of small scale dissipation in the above
calculation, but in the time-average, small scale dissipation simply
decreases the net generation rate by a calculable amount.  This will
be discussed further in the context of the numerical simulations.

\subsection{Dimensional estimates of halting scale}
\label{simpest:ssc}

Using \eqref{TA:eq}, setting $T_A = r^{-1}$ and assuming
$g=\epsilon$ gives an estimate for the halting scale
\begin{equation}
k_r \sim \left(r^3g^{-1}\right)^{{1 / [3(\gamma-1)]}}.
\label{krdim:eq}
\end{equation}
For our three special cases this yields
\begin{subequations}
\begin{align}
\textrm{TDV:} \qquad &k_r \sim \left(r^3g^{-1}\right)^{1/2} \\
\textrm{SQG:} \qquad &k_r \sim \left(r^3g^{-1}\right)^{1/3} \\
\textrm{LQG:} \qquad &k_r \sim \left(r^3g^{-1}\right)^{1/8}
\end{align}
\end{subequations}
where $g$ is the injection rate of the appropriate upscale-cascading
invariant.

The expression \eqref{krdim:eq} is demanded by dimensional
consistency, and so gives no indication of the value of any numerical
prefactor.  One simple way to derive such a prefactor is to note that
the integral of the energy spectrum \eqref{Aspec:eq} must be
consistent with the total energy as given by \eqref{Ar2:eq}. That is,
\begin{equation}
\int_{k_r}^\infty \ca{C}\epsilon^{2/3}k^{-\gamma} \, dk = \frac{g}{2r},
\label{krint:eq}
\end{equation}
where we have assumed that the energy content in the larger scales
($k<k_r$) is negligible.  For $\gamma > 1 $ this gives
\begin{equation}
k_r \approx \left(\frac{2\ca{C}}{\gamma-1} \right)^{1/(\gamma-1)}
\left(r^3g^{-1}\right)^{{1 / [3(\gamma-1)]}}.
\label{krgen1:eq}
\end{equation}
which is identical to \eqref{krdim:eq}, except for the numerical
coefficient involving the Kolmogorov constant. For TDV, this gives
\begin{equation}
k_r \approx \left({3\ca{C}} \right)^{3/2}
\left(r^3g^{-1}\right)^{1/2}.
\label{krtdv1:eq}
\end{equation}
\citet{danilov01a} independently proposed an identical scaling.  This
coefficient is accurate only to the extent that the energy falls off
very rapidly for wavenumbers smaller than $k_r$.

\subsection{Spectral shapes}
\label{univ:ssc}

By making one additional assumption, we can derive an expression for
the spectral shape that is putatively valid in the inertial range and
in the range where friction begins to dominate.  The peak of this
theoretical spectrum, derived in a manner similar to that of
\citet{Leith68a} and \citet{Lilly72}, yields a better estimate for the
stopping-scale prefactor than \eqref{krgen1:eq}. However, we only
expect the theory to be relevant in parameter regimes such that the
underlying inertial range is present and free of intermittency, and
make no claim that the derived spectral shape is universal.

The spectral budget equation for generalised energy (enstrophy) in
mode $\kb$ for $\alpha>0$ ($\alpha<0$) is
\begin{equation}
\pp{A_{\kb}}{t} = \epsilon_{\kb} - 2rA_{\kb} + g\delta(k-k_f),
\label{A:eq}
\end{equation}
where $A_{\kb} = |\psi_{\kb}|^2/2$ (or
$\mathrm{Re}\left[\xi_{\kb}^{\star}\psi_{\kb}\right]/2$),
$\epsilon_{\kb}$ is the transfer rate at $\kb$ arising from the real
part of the product of $-\psi_{\kb}^{\star}$ and the spectral Jacobian
term, and $g$ is the generalised energy (enstrophy) injection,
localized at isotropic wavenumber $k_f$.  We assume a steady state, so
that the time derivative vanishes, and isotropy so that $\kb$ can be
replaced by $k$.  If the source term is at some high wavenumber $k_f
\gg k_r$ then we may expect an inverse cascade initially unfettered by
friction.  At smaller wavenumbers, the drag will remove energy from
the flow, and, without further approximation, the spectral transfer is
governed by the budget equation
\begin{equation}
\dd{\epsilon}{k} = 2r\ca{A}(k).
\label{epsr:eq}
\end{equation}
That is, the energy flux is reduced by the frictional loss.

We now assume that the spectrum \eqref{Aspec:eq} can be substituted
for $\ca{A}(k)$, but let $\epsilon = \epsilon(k)$ therein. This is a
perturbative solution for small drag-induced deviations from the
inertial range flux.  Thus we expect the approximation to be valid in
the spectral regions where drag is small and up to the scales where it
begins to dominate, but not at the frictionally dominated, largest
scales.  In this approximation \eqref{epsr:eq} becomes
\begin{equation}
\dd{\epsilon}{k} = 2r\ca{C}\epsilon^{2/3}k^{-\gamma}
\label{epsr2:eq}
\end{equation}
which readily yields
\begin{equation}
\epsilon(k) = \left[\frac{2\ca{C}}{3(1-\gamma)}rk^{1-\gamma} +
c\right]^3
\label{epsr3:eq}
\end{equation}
where $c$ is the integration constant.  The transfer rate must be
equal to the generation rate at the forcing scale, \ie
\begin{equation}
\epsilon(k_f) = g,
\label{bc:eq}
\end{equation}
which we use to solve for the integration constant in
\eqref{epsr3:eq}.  Using the solution for $c$, the flux is
\begin{equation}
\epsilon(k) =
g~\left[1 - (k_c/k)^{\gamma-1} + (k_c/k_f)^{\gamma-1} \right]^3,
\label{epsr4:eq}
\end{equation}
where
\begin{equation}
k_c =\left[\frac{2\ca{C}}{3(\gamma-1)}\right]^{1/(\gamma-1)}
\left(r^3g^{-1}\right)^{1/[3(\gamma-1)]}.
\label{kc:eq}
\end{equation}
Substituting this into \eqref{Aspec:eq} gives us an estimate for the
spectrum of the invariant in the presence of linear drag
\begin{equation}
\ca{A}(k) = \ca{C}g^{2/3}k^{-\gamma}
\left[1 - (k_c/k)^{(\gamma-1)} + (k_c/k_f)^{\gamma-1}\right]^2.
\label{Arspec:eq}
\end{equation}
For TDV, the expression \eqref{Arspec:eq} in the limit $k_f
\rightarrow \infty$ was derived by \citet{Lilly72}.  The reader
may verify that the expression for $\ca{A}(k)$ satisfies
\begin{equation}
\int_{k_c}^{\infty}\ca{A}(k)\, dk = \frac{g}{2r}
\label{Aint:eq}
\end{equation}
exactly.

For $k_f\gg k_c$, the third term in brackets is small and can be
neglected (but we shall retain it in our predictions for the simulated
spectra).  For wavenumbers $k\gg k_c$, the second term in square
brackets in \eqref{Arspec:eq} is also negligible and the spectrum
reduces to \eqref{Aspec:eq}. For sufficiently small wavenumbers,
$\epsilon(k)$ given by \eqref{epsr4:eq} is not positive. The critical
wavenumber at which the spectral flux vanishes is $k=k_c$ (assuming
$k_f\rightarrow\infty$), where $k_c$ is given by \eqref{kc:eq}. It is
unphysical for $\ca{A}(k)$ to vanish identically at some finite
$k$. The weakness of this derivation is just the use of the inertial
range spectrum \eqref{Aspec:eq} in \eqref{epsr:eq}, an assumption we
expect to be valid only if frictional effects are small perturbations
on the inertial flow, as mentioned above.

The predicted spectral shapes for the three cases we consider here are
(assuming $k_f\rightarrow\infty$)
\begin{subequations}
\begin{alignat}{3}
&\textrm{TDV:}
\qquad \ca{E}(k) = \ca{C}g^{2/3}k^{-5/3}
\left[1-(k_c/k)^{2/3}\right]^2,
\qquad &k_c& = \ca{C}^{3/2}\left(r^3g^{-1}\right)^{1/2},
 \label{Erspec2d:eq}\\
&\textrm{SQG:}
\qquad \ca{E}(k) = \ca{C}g^{2/3}k^{-2}
\left[1-(k_c/k)\right]^2,
\qquad &k_c& = \left(2\ca{C}/3\right)
 \left(r^3g^{-1}\right)^{1/3},\label{Erspecsqg:eq}\\
&\textrm{LQG:}
\qquad \ca{Z}(k) = \ca{C}g^{2/3}k^{-11/3}
\left[1-(k_c/k)^{8/3}\right]^2,
\qquad &k_c& = \left(\ca{C}/4\right)^{3/8}
 \left(r^3g^{-1}\right)^{1/8},\label{Erspeclqg:eq}
\end{alignat}
\end{subequations}
where in each case the spectrum is valid for all $k>k_c$, and $\ca{C}$
and $g$ are the Kolmogorov constant and injection rate, respectively,
for the particular dynamics at hand.

A quantitative estimate of the drag-induced halting scale can be made
by calculating the location of the peak wavenumber $k_0 = k_r$ by
setting
\begin{equation}
\left.\dd{\ca{A}(k)}{k}\right|_{k=k_r} = 0.
\label{dAdk0:eq}
\end{equation}
This yields
\begin{equation}
k_r = \Gamma\left(r^3g^{-1}\right)^{1/[3(\gamma-1)]},~~~
\Gamma = \left[\frac{2\ca{C}(3\gamma-2)}{3\gamma(\gamma-1)}
\right]^{1/(\gamma-1)}.
\label{krAtrue:eq}
\end{equation}
For our three cases, the drag-induced stopping scales expected
are
\begin{subequations}
\begin{align}
\textrm{TDV:}
\qquad & k_r = \left(9\ca{C}/5\right)^{3/2}
\left(r^3g^{-1}\right)^{1/2},\label{kr2d:eq}\\
\textrm{SQG:}
\qquad & k_r = \left(4\ca{C}/3\right)
\left(r^3g^{-1}\right)^{1/3},\label{krsqg:eq}\\
\textrm{LQG:}
\qquad & k_r = \left(27\ca{C}/44\right)^{3/8}
\left(r^3g^{-1}\right)^{1/8}.\label{krlqg:eq}
\end{align}
\end{subequations}

As in \eqref{krtdv1:eq}, the prefactor $\Gamma$ in each case is not
necessarily an order unity quantity. In the case of energy in TDV
turbulence, for example, using the estimate $\ca{C} = 5.8$
\citep{Maltrud91} yields $\Gamma = (9\ca{C}/5)^{3/2} = 33.7$. Using
\eqref{krtdv1:eq} gives $72.6$, the larger number being consistent with the
assumption that energy is strictly zero for wavenumbers smaller than $k_r$.
Our numerical results (discussed in \S \ref{numerics:sec}) are consistent
with these estimates; that is, the strictly dimensional estimate $k_r =
(r^3/g)^{1/2}$ is more than an order of magnitude smaller than the actual
wavenumber at which the cascade is halted.  The large value of $\Gamma$
indicates that, in the presence of drag, TDV turbulence is more inefficient
at moving energy to larger scales than one might expect from dimensional
estimates alone.


\section{Mean vorticity gradients:  the $\beta$-effect}
\label{beta:sec}

In two-dimensional vorticity dynamics, a meridional gradient of
planetary vorticity may be represented by the $\beta$-effect. In
SQG, a meridional gradient of temperature gives rise to a similar
term in the equation of motion. Including such effects in
\eqref{gen2dr:eq}, we write
\begin{equation}
\pp{\xi}{t} + J(\psi,\xi) +\beta\pp{\psi}{x} = F-r\xi,
\label{gen2drb:eq}
\end{equation}
where we use the symbol $\beta$ for the constant mean gradient of
$\xi$ with units $\beta \sim [L^{1-\alpha}T^{-1}]$.  Because $\beta$
introduces anisotropy into the dynamics \citep{Rhines75,Vallis93a},
the halting scale, if present, will be anisotropic as well. However,
the dual presence of $\beta$ and $r$ introduces complexities to the
dynamics which, even for TDV, have only recently come to be
appreciated
\citep{chekhlov96,manfroi99,smithwaleffe99,huang01,galperin01,danilov01b}. The
mean vorticity gradient admits a possible tri-partite balance in the
equation of motion between non-linear advection, advection of the mean
gradient (and thereby possible wave propagation) and dissipation.

We direct the discussion here to the two cases with $\alpha>0$: TDV
and SQG.  The $\beta$ term as defined in \eqref{gen2drb:eq} physically
represents the Coriolis gradient in TDV, and a mean temperature
gradient in SQG, but has no obvious physical interpretation for
LQG\footnote{Keeping the same definition of $\beta$ as in TDV, we
obtain a term of the form $\beta\partial{\xi}/\partial{x}$ added to
the left hand side of \eqref{lqg3:eq}. The dispersion relation is the
same as that for equivalent barotropic Rossby waves at large scales,
namely \(\omega_{R,LQG} = -k_x\beta/\lambda^2.\) The frequency goes
linearly to zero rather than diverging at large scale, hence it is not
clear that $\beta$ can halt the cascade in this case.}.

\subsection{Waves, turbulence and drag}

The first significant scale for these flows is determined by the
balance between non-linear advection and advection of the mean
gradient.  Ignoring drag for the moment, Rossby waves in the
generalised two-dimensional formalism will have the dispersion
relation
\begin{equation}
\omega_R = -\frac{k_x\beta}{k^{\alpha}}.
\label{rossby:eq}
\end{equation}
Following \citet*{Vallis93a} we set this equal to the eddy strain rate
(the inverse of \ref{TA:eq}) to find the boundary between
$\beta$-dominated and isotropic flow,
\begin{equation}
k(\theta) = k_{\beta}~(\cos \theta)^{1/(\alpha+\gamma-2)}
\label{wtbound:eq}
\end{equation}
where
\begin{equation}
k_{\beta} = (\ca{C}^{-1/2}\beta\epsilon^{-1/3})
^{1/(\alpha+\gamma-2)}.
\label{kbeta1:eq}
\end{equation}
Note that in the case of energy in TDV dynamics
($\alpha=2$ and $\gamma=5/3$), expression \eqref{kbeta1:eq} becomes
\begin{equation}
k_{\beta} \simeq \left(\frac{\beta^3}{\ca{C}^{3/2}\epsilon}\right)^{1/5},
\label{kbetavm:eq}
\end{equation}
a form similar to that of \citet*{Vallis93a}, apart from the
nondimensional factors in the denominator (which we keep for later
use).

The relation \eqref{wtbound:eq} specifies a parametric relationship
between the modulus of wavenumber $k$ and the wave direction $\theta =
\arctan(k_y/k_x)$ having the shape of a dumbbell in the $(k_x,k_y)$
plane and having a maximum at $\theta=0$ (along the $k_x$ axis).  Thus
for $\beta$ such that $k_{\beta}\ll k_f$, we expect an isotropic
inverse cascade up to $k_{\beta}$, and anisotropic flow at wavenumbers
below $k_{\beta}$, rather than a simple halting of the inverse cascade
at $k_{\beta}$.  Because weakly non-linear wave interaction typical of
anisotropic flows is not as efficient as isotropic turbulence at
cascading the invariant, we expect this boundary to be significant
also in the spectral distribution of energy.  Along the $k_y$ axis,
\eqref{wtbound:eq} gives $k(\theta=\pi/2)=0$.  There is evidence
\citep{Vallis93b,chekhlov96,smithwaleffe99, huang01} that energy (in
TDV flow on the torus or sphere) will continue to cascade to larger
scale past $k_{\beta}$, but with most of the energy concentrated along
the $k_y$ axis.

When drag is present, and the drag coefficient $r$ is sufficiently
large, a balance can be reached between dissipation and non-linear
advection at scales smaller than those at which $\beta$ becomes
important, and the cascade is then halted isotropically.  A critical
value exists for $r$, separating the isotropic steady state flow from
one in which anisotropy is allowed to develop.  This critical value of
$r$ is found by setting $k_r = k_{\beta}$, using \eqref{krAtrue:eq}
and \eqref{kbeta1:eq}, yielding
\begin{equation}
r_c = \Gamma^{1-\gamma}g^{(1-\alpha)/[3(2-\gamma-\alpha)]}
(\ca{C}^{-1/2}\beta)^{(1-\gamma)/(2-\gamma-\alpha)}
\label{rcrit:eq}
\end{equation}
In the TDV case with $\alpha=2$ and $\gamma=5/3$ this expression gives
\begin{equation}
r_{c,\alpha=2} = \frac{5}{9\ca{C}}(\beta^2g)^{1/5}.
\label{rcrit2d:eq}
\end{equation}
In the rest of this discussion, we will assume that $r<r_c$, so that
an anisotropic, large-scale spectrum should ensue.  It then remains to
determine the largest scale to which this anisotropic spectrum
extends, and expecting this scale to depend on both $\beta$ and $r$,
we denote it $k_{\beta,r}$.

\subsection{Estimates of the halting scale with $\beta$ and drag present}

For clarity of presentation we focus on TDV, and then discuss how it
may be adapted to generalised two-dimensional turbulence.  In the
small drag limit ($r \ll r_c$) we seek an estimate for the halting
scale $k_{\beta,r}$, that is for the scale at which the maximum energy
is found.  Because in anisotropic turbulence most of the energy goes
into zonal jets it is not surprising that this estimate will be
particularly apt as an estimate for the separation scale of the jets
themselves.  Our argument has two components: (i) the forcing and
dissipation set the rms velocity; (ii) using this in conjunction with
the $\beta$-effect provides an estimate of the halting scale.  We
shall see that despite the jet scale being the most energetic scale,
it is the inviscid scale \eqref{kbeta1:eq} that characterises the
meridional transport of passive tracer.

If one assumes that, at scales where $\beta$ is important but where
drag is not yet significant ($k_{\beta,r}<k<k_{\beta}$), the spectral
slope is a function only of $\beta$ and wavenumber, then dimensional
analysis leads to the expression
\begin{equation}
\label{rossbyspecB:eq}
\ca{E}_\beta(k) = \ca{C}_{\beta} \beta^2 k^{-5},
\end{equation}
as pointed out by \citet{Rhines75}.  Obviously, though, the spectrum
cannot be isotropic, since it is derived specifically in the region
where anisotropy is present.  Our simulations, as well as those of
\citet{galperin01} and others cited above, demonstrate that at scales
$k_{\beta,r}<k<k_{\beta}$, energy is concentrated along the $k_x=0$
axis, but does roughly follow a $k^{-5}$ power law.
\citet{chekhlov96} point out that energy continues to follow a
$k^{-5/3}$ power law for all wavenumbers $k_{\beta,r}<k<k_{\beta}$
with $k_x \neq 0$, but that implies that modes with $k_x \neq 0$
contribute negligibly to the total spectrum at large scale.

\mbox{} \citet{smithwaleffe99} discuss the subtleties of the transfer from the
two-dimensional isotropic cascade to the one-dimensional ``cascade''
along the zonal axis.  If the energy is concentrated on the $k_x=0$
axis, precisely where the $\beta$ term in \eqref{gen2drb:eq} vanishes,
then the dependence on $\beta$ of \eqref{rossbyspecB:eq} is
unexpected.  A dependence on $\beta$ might arise if resonant nonlinear
triad interactions among Rossby waves were responsible for the
transfer to zonal flow, but these are unable to transfer energy from
isotropic two-dimensional motions to modes with $k_x =
0$. \citet{smithwaleffe99} argue that it is necessary to invoke
quartic resonances or near resonances to explain these transfers and
still maintain a $\beta$ dependence in the
spectrum. \citet{chekhlov96} make similar points.  \citet{danilov01b}
point out that in fact the zonal energy is not distributed as a
uniform spectrum, but rather as a series of peaks whose amplitudes
increase roughly like $k^{-5}$, and provide a preliminary explanation
for the observed wavenumber dependence.  Given all of the above, we
will take the appearance of such a $-5$ spectrum (or wavenumber
dependence of the spectral peaks) as an empirical observation with a
rather weak theoretical or phenomenological justification.

Thus, we will suppose an isotropic $k^{-5/3}$ spectra at scales
$k_{\beta}<k<k_f$, an anisotropic $k^{-5}$ spectral distribution, with most
energy concentrated along the zonal axis, at scales
$k_{\beta,r}<k<k_{\beta}$, and negligible energy at $k<k_{\beta,r}$. That
is,
\begin{equation}
\ca{E}(k) \simeq
\left\{\begin{array}{ll}
    \ca{C}\epsilon^{2/3}k^{-5/3},  &\mbox{$k_{\beta}<k<k_f$}   \\
    \ca{C}_{\beta} \beta^2 k^{-5}, &\mbox{$k_{\beta,r}<k<k_{\beta}$,
                                          ~$k_x=0$}\\
    0,                             &\mbox{$k<k_{\beta,r}$}.
\end{array}
\right.
\label{betaspectot:eq}
\end{equation}

We use this phenomenological picture to determine $k_{\beta,r}$.  Note
that even in the presence of $\beta$, \eqref{Atot:eq} remains valid,
so that the total energy is $E = g/2r$. This constraint determines the
lowest wavenumber which the inverse cascade may reach, just as in
\eqref{krgen1:eq}.  Specifically, using \eqref{betaspectot:eq} we set
\begin{equation}
\int_{k_{\beta,r}}^{\infty} \ca{E}(k) \, dk \simeq \frac{g}{2r}.
\label{intAros:eq}
\end{equation}
Because the spectra are so steep in the Rossby wave regime, the
integral is dominated by contributions from its peak at the low
wavenumber end, giving
\begin{equation}
k_{\beta,r} \simeq \left(\frac{\ca{C}_{\beta}\beta^2r}{2g}\right)^{1/4}.
\label{kbr:eq}
\end{equation}

A simpler derivation results if we begin with the
familiar expression for the stopping scale \citep{Rhines75}
\begin{equation}
k_{\beta,r} = \left(\beta U_{\rms}^{-1}\right)^{1/2}
\label{rh75:eq}
\end{equation}
But from \eqref{Ar2:eq} $U_{\rms} = (g/r)^{1/2}$ (where we have also
assumed that $U_{\rms}\gg V_{\rms}$), and substituting this into
\eqref{rh75:eq} we obtain \eqref{kbr:eq}, albeit without the numerical
constant. Note that \eqref{rh75:eq} is an estimate for a marginally
barotropically unstable jet.

\mbox{} \citet*{manfroi99} predict a jet-separation scale pro\-por\-tion\-al
to $r^{-1/3}$, or equi\-va\-lent\-ly, $k_{\beta,r} \propto r^{1/3}$.
Distinguishing between $r^{-1/3}$ and $r^{-1/4}$ (the latter from our
\eqref{kbr:eq}) is difficult numerically due both to the small difference
in exponent and to the slow equilibration of the flow in this parameter
regime.

We can formally derive an anisotropic wave spectrum in the general
case assuming a balance between the wave frequency \eqref{rossby:eq}
and the inverse of the eddy timescale of \eqref{TA:eq}. The wave
spectrum in this case is
\begin{equation}
\ca{A}_{\beta}(k) = \ca{C}_{\beta} \beta^2 k^{4-2\alpha-3\gamma},
\label{rossbyspecA:eq}
\end{equation}
which reduces to \eqref{rossbyspecB:eq} in the inverse cascade of
TDV. In SQG we find a $k^{-4}$ spectrum in the Rossby
wave regime, also significantly steeper than the isotropic spectrum of
$k^{-2}$ relevant in that case.

The result analogous to \eqref{kbr:eq} for the general case is
\begin{equation}
k_{\beta,r} \simeq \left[\frac{2\ca{C}_{\beta}\beta^2r}
{(2\alpha+3\gamma-5)g}\right]^{1/(2\alpha+3\gamma-5)}
\label{kbrA:eq}
\end{equation}
The simple argument using the analog of \eqref{rh75:eq} is not always
possible in the generalised case, because the integral constraint
corresponding to \eqref{Atot:eq} does not always give a value for the
rms velocity. For example, in SQG, the generalised energy
$E=\overline{\psi\xi}/2$ is not a velocity squared. Rather, it is the
generalised enstrophy that has units of velocity squared, but its value
is not constrained solely by the forcing and the drag.

Which of the scales \eqref{kbetavm:eq} and \eqref{kbr:eq} (or
\eqref{kbeta1:eq} and \eqref{kbrA:eq} in the general case) is relevant
for the meridional transport of a passive tracer? The scale
\eqref{kbr:eq} determines a jet scale, with vanishing meridional
velocity. \citet{Bartello91} demonstrate that at the jet scale, zonal
advection so dominates meridional advection that little meridional
mixing occurs.  We thus expect the scale \eqref{kbetavm:eq}, the
largest scale of \textit{isotropic} turbulence, to act as a mixing
length, and explore this further in \S \ref{numerics:sec}.
\citet{Held96} utilized the scaling \eqref{kbetavm:eq}, without the
numerical prefactor, as a mixing length in their theory for turbulent
quasi-geostrophic heat fluxes.

\section{Numerical examinations of diffusivity scalings}
\label{numerics:sec}

We now describe sets of numerical simulations designed to test the
diffusivity predictions for TDV, SQG and LQG with linear drag.  We
also test TDV with linear drag and $\beta$.  Mixing lengths for each
case have been derived in the previous two sections, but velocity
scales have not --- both are required for estimates of diffusivity
given by \eqref{diff:eq}.  It is difficult to generalise the
calculation of velocity variance among the various cases, so we
consider this detail, and the predictions for the diffusivities
themselves, in the context of each case separately.

All simulations were performed using a two-dimensional de-aliased
spectral model with $512^2$ equivalent horizontal gridpoints
($k_{\mathrm{max}}/k_{\mathrm{min}} = 255$). A leap-frog timestep is
used to advance the solution, and a weak Robert filter suppresses the
computational mode.  In each case, a passive tracer is advected along
with the calculated flow.  The flow is forced with an isotropic
forcing at high wavenumber, typically about $k_f/k_{\mathrm{min}} =
160$.  The magnitude of the generation rate is fixed at $g=1$, less
some due to small-scale dissipation (see appendix B for details of the
forcing function).

Small scale variance in both the generalised vorticity field (\ie
enstrophy) and in the tracer field are dissipated with a highly
scale-selective exponential cutoff filter (\ref{fec:eq} and
\ref{a:eq}) with vanishing dissipation below a cutoff wavenumber
$k_{\mathrm{cut}}$.  For all simulations presented here,
$k_{\mathrm{cut}}/k_{\mathrm{min}}=165$ for the advecting flow and
$k_{\mathrm{cut}}/k_{\mathrm{min}}=1$ for the tracer.  The exponent
$s$ in \eqref{fec:eq} is $s=8$ for the advecting flow and $s=2$ for
the tracer.  For the advecting flow, $k_{\mathrm{cut}}$ is just larger
than the forcing scale and is the scale at which we begin to allow
small-scale dissipation --- traditional hyperviscosity, by contrast,
dissipates at all scales. The dissipation level does not become
significant until near the maximum wavenumber in the computational
domain.  Thus, while we focus here on maintaining as wide an inverse
cascade range as possible, we have allowed a reasonable direct cascade
for wavenumbers $k>k_f$.  The small-scale dissipation, despite its
presence only at $k>k_f$, removes some of the upscale-cascading
invariant.  We calculate this loss and use it to define a
$g_{\mathrm{eff}} < g = 1$ ($g_{\mathrm{eff}} \simeq .5$ for all the
simulations --- see appendix B).

In each of the following sections, we state values as
non-dimensionalized using a length given by the inverse forcing
scale $L = k_f^{-1}$ and a time given by a combination of the
generation rate $g$ and forcing scale.  The former has units
\[
g \sim \left[L^{3(\gamma-1)}T^{-3}\right],
\]
so that a nondimensional measure of the time is
\begin{equation}
T = \left[g k_f^{3(\gamma-1)}\right]^{-1/3}.
\label{Tdim:eq}
\end{equation}
All values with a $\star$ superscript have been
non-dimensionalized by these scales.

For each dynamical case considered, a series of simulations was
performed.   Each simulation described is run several eddy
turnover times after a statistically steady state is achieved;
averages are taken over the equilibrated phase of a given run.
Each simulation was performed with a different drag coefficient,
the values of which were chosen such that the halting scales were
large compared to the forcing scale, but not too close to the
domain scale.

\subsection{Two-dimensional vorticity flow with linear drag}
\label{tdv_exp:ssc}

We begin with the straightforward case of two-dimensional
vorticity dynamics dissipated by linear drag, namely equations
\eqref{gen2dr:eq} and \eqref{coupling:eq} with $\alpha=2$.  In
this case $E$ cascades upscale with spectral slope
$\gamma=5/3$.  Five simulations were performed in which the drag
was set to
\[
r^{\star} = r~\left(gk_f^2\right)^{-1/3}
=(0.86, 1.5, 2.4, 3.7, 6.0)\times 10^{-2}.
\]

We predict a spectral shape given by \eqref{Erspec2d:eq} and a
spectral peak given by \eqref{kr2d:eq}.  Since the total energy,
according to \eqref{Ar2:eq}, is
$E=g/2r=(U_{\rms}^2+V_{\rms}^2)/2\simeq V_{\rms}^2$ (assuming
isotropy), the integrated (rms) meridional velocity $V_{\rms}$ is
\begin{equation}
V_{\rms} = \left(\frac{g}{2r}\right)^{1/2}.
\label{V2d:eq}
\end{equation}
Using \eqref{diff:eq}, \eqref{kr2d:eq} and \eqref{V2d:eq}, we
estimate the eddy diffusivity, or integrated tracer flux, to be
\begin{equation}
D \simeq
\left[\left(\frac{5}{9\ca{C}}\right)^3\frac{1}{2}\right]^{1/2}gr^{-2}.
\label{diff2dr:eq}
\end{equation}
We use $\ca{C}=6$ in numerical estimates.

Figure \ref{tdv:fig} shows the spectra of the energy, tracer variance
and tracer flux, as well as the stopping scales, rms velocities and
diffusivities for each simulation, along with their respective
predictions.  The halting scale $k_0$ is predicted by \eqref{kr2d:eq}
with some accuracy.  [A similar result is inferred from the soap-film
experiments of \citet{Rivera00}.]  However, the amplitude of the
predicted spectrum \eqref{Erspec2d:eq} falls off much more rapidly
than the experimental spectrum at small wavenumbers; the predicted
spectrum evidently fails when drag is dominant.  Consistently, the
largest overall discrepancy between theory and simulation occurs for
the simulation with the highest drag, in which case the magnitude of
the predicted spectrum falls well below that obtained numerically.
Neither fact is surprising, since the theoretical spectrum is derived
under the assumption of significant separation between the drag and
forcing scales, and is formally valid only where deviations from
constant $\epsilon$ are small, which is not true at small wavenumbers.
Both the tracer variance and flux spectra have slopes close to $-5/3$,
as predicted by \eqref{PE:eq} and \eqref{Dspec:eq}, respectively.

The physical-space streamfunction and tracer fields for the run with
the lowest drag are shown in the upper two panels of figure
\ref{phys:fig}. Streamfunction, rather than vorticity, is plotted
since the latter is dominated by noise from the random forcing field.
The large-scale flow represented in the streamfunction is nearly
structureless, but with undulations whose scale represents that of the
spectral peak. The absence of coherent vortices in the flow is
consistent with the good fit of the theory for the spectrum in figure
\ref{tdv:fig}(a).  The tracer field, plotted with its mean gradient,
is composed of meridionally-oriented plumes, with a characteristic
scale similar to that of the streamfunction.  (The presence of fronts,
and ``sheets'' in the passive tracer field has been investigated by
\citet{constantin94}. %, and more recently by \citet{celani01}.
% Personally, I'd also remove the constantin reference. 

\subsection{Surface quasi-geostrophic flow}
\label{sqg_exp:ssc}

We now present simulations of SQG flow with linear drag, \ie equations
\eqref{gen2dr:eq} and \eqref{coupling:eq} with $\alpha=1$.  Since
$\alpha>0$, the upscale cascading invariant has a spectral slope
$\gamma=2$.  Five simulations were performed, differing only in the
magnitude of drag, whose value was set to
\[
r^{\star} = r\left(gk_f^3\right)^{-1/3}
=(0.18,0.35,0.72,1.4,2.9)\times 10^{-2},
\]
respectively.

The spectral shape and peak are predicted by \eqref{Erspecsqg:eq} and
\eqref{krsqg:eq}, respectively.  Estimating the velocity in this case
is less straightforward.  As defined in section \ref{spectra:sec}, for
$\alpha=1$, the invariants of the flow have units $E\sim[L^3T^{-2}]$
and $Z\sim[L^2T^{-2}]$, thus in this case $Z$ has the dimensions of
energy (physically it is the temperature variance, or available
potential energy).  In SQG then, the rms velocity of the flow is
related to the generalised enstrophy as $Z=V_{\rms}^2$ (assuming
isotropy).  Ignoring drag for the moment, in the inverse cascade range
$Z$ has the spectrum
\begin{equation}
\ca{Z}(k) = \ca{C}_E\epsilon_E^{2/3}k^{-1}.
\label{Zsqg:eq}
\end{equation}
The integral of this spectrum diverges logarithmically, and we must
explicitly include the forcing scale
\begin{equation}
Z \simeq \int_{k_c}^{k_f}\ca{Z}(k)\, dk,
\label{Zsqg2:eq}
\end{equation}
where $k_c$ is as given by \eqref{Erspecsqg:eq}.  We thus obtain our
estimate of the rms meridional velocity scale
\begin{equation}
V_{\rms} = \left[\ca{C}_Eg^{2/3}\ln\left(k_f/k_c\right)\right]^{1/2},
\label{Vsqg:eq}
\end{equation}
where we have taken $g=\epsilon_E$.

Using \eqref{diff:eq}, \eqref{krsqg:eq} and \eqref{Vsqg:eq}, we
estimate the eddy diffusivity, or integrated tracer flux, to be
\begin{equation}
D \simeq
\left(g^{2/3}r^{-1}\right)
\left[\frac{9}{16\ca{C}_E}\ln\left(k_f/k_c\right)\right]^{1/2}.
\label{diffsqg:eq}
\end{equation}

The spectra of $\ca{E}(k)$, the tracer variance and tracer flux, as
well as the stopping scales, rms velocities and diffusivities, and
their respective theoretical predictions, are plotted in figure
\ref{sqg:fig}.  The predictions for the spectral shapes in the upper
left panel are noticeably less accurate than were the analogous curves
for TDV in the previous section.  In particular, the spectra are more
shallow than the expected $-2$ slope, even at scales which fall well
inside the inverse cascade inertial range.  One might attribute this
deviation from the expected slope to intermittency caused by the
formation and persistence of coherent vortices in the flow.  Indeed,
we see evidence for such structures, but in this case one would expect
the spectral slope to be steeper than its inertia range theory
prediction, when in fact it is somewhat more shallow.  We offer no
explanation for this behavior presently.

We estimate the Kolmogorov constant such that the closest possible fit
was achieved for all five cases, resulting in the value $\ca{C} =
7.9$.  Despite the poor fits, the simulated stopping scales are fairly
well predicted by the peaks of the theoretical spectral shapes.  Both
tracer variance and flux have the expected slopes $-2$.  Furthermore,
despite the problematic dependence of the rms velocity on the forcing
scale in \eqref{Vsqg:eq}, the predictions made thereby are relatively
accurate.  The kinks in the theoretical curves for the velocities and
diffusivities are due to the fact that the effective generation rate
$g_{\mathrm{eff}}$ changed somewhat from one run to the next (see
appendix B).  The diffusivity scales as anticipated, but is
overestimated in magnitude by about a factor of 2.

The physical-space streamfunction and tracer fields for the run with
the lowest drag are shown in the middle two panels of figure
\ref{phys:fig}.  The structures of the flow and tracer are similar to
those for TDV, but with a higher degree of localisation in the
streamfunction field, due to the presence of some small-scale coherent
vortices in the flow.  We nevertheless see patterns in the tracer
field similar to those found for TDV flow in the previous subsection.

\subsection{Large-scale quasi-geostrophic flow}
\label{lqg_exp:ssc}

Here we consider the case \eqref{gen2dr:eq} with $\alpha=-2$, termed
LQG.  We actually integrate the shallow water quasi-geostrophic
equation
\begin{equation}
\pp{q}{t} + J(\psi,q) = F - rq, ~~~ q = (\del^2 - \lambda^2)\psi
\label{qgswr:eq}
\end{equation}
with $\lambda/k_{\mathrm{min}}=70$, or $\lambda/k_f = 0.44$ so
that at small wavenumbers the dynamics should be described by
\eqref{gen2dr:eq} with $\alpha=-2$.  By integrating
\eqref{qgswr:eq} one may see the transition from QG to LQG near
scales just larger than the deformation scale (a similar approach
was used by \citet{kukharkin95}).

In the LQG regime ($k\ll\lambda$), the generalised energy and
enstrophy are in fact the kinetic energy (KE) and available potential
energy (APE), respectively.  One can see this by using the variable
change \eqref{lqgvar:eq} in the general expressions for the
invariants \eqref{invar:eq}, which yields
\begin{equation}
E_G = \frac{1}{2}\overline{\Psi\xi} =
\frac{1}{2}\sum_{\kb}k^2|\psi_{\kb}|^2 = KE,
\label{lqgE:eq}
\end{equation}
and
\begin{equation}
Z_G = \frac{1}{2}\overline{\xi^2} =
\frac{1}{2}\sum_{\kb}|\psi_{\kb}|^2 = APE/\lambda^2.
\label{lqgZ:eq}
\end{equation}

Since $\alpha < 0$, the generalised enstrophy, or APE, cascades
upscale, while the generalised energy, or KE, cascades toward small
scales.  Because we have chosen the deformation wavenumber
$\lambda<k_f$, we expect a transition from the TDV-like inverse
cascade at scales such that $k\gg\lambda$ to the LQG inverse cascade
at scales $k\ll\lambda$.  Furthermore, there will be no direct cascade
range for LQG in this arrangement, since the forcing scale is at
smaller scale than the transition region.

In the LQG regime, from \eqref{gam:eq}, $\alpha = -2$ implies that
for $Z$ (the APE) the spectral slope is $\gamma=11/3$.  But since
$\ca{Z}_{\mathrm{LQG}}(k) = |\psi^2|$, the spectral slope of the
kinetic energy $\ca{E}_{\mathrm{LQG}}(k) = k^2|\psi^2|$ is
$-5/3$, just as it is for $k^2|\psi^2|$ in the TDV regime.  Thus
we expect \textit{no change in slope} of $k^2|\psi^2|$ as we
transit from TDV to LQG, but possibly a kink in the spectrum
due to a possible change in the Kolmogorov constant.

In figure \ref{lqg_spec:fig} we have plotted the spectrum of
$\ca{Z}(k)$ for the run with the smallest drag, along with the
theoretical prediction described below.  Notice that at scales just
larger than the deformation scale (indicated by the arrow), the
spectrum slumps but maintains a similar slope, rising slightly near
the spectral peak.  Because of the limited extent of the inertial
range, an accurate determination of the appropriate Kolmogorov
constant is not possible, and we use the value from TDV, $\ca{C}= 6$.

We expect the spectrum of generalised enstrophy to follow
\eqref{Erspeclqg:eq} and the peak of this spectrum to be given by
\eqref{krlqg:eq}.  These expressions refer to the LQG formalism
as derived in \S \ref{spectra:sec}, in which $g$ is the
generation rate of the generalised enstrophy.  Dimensionally $g
\sim [L^8\tau^{-3}]$, where $\tau$ is the rescaled time
coordinate (see equation \ref{tau:eq}).  Thus in terms of
unscaled time, $\tilde{g} = \lambda^6g$, hence
\begin{equation}
k_c = \left(\ca{C}/4\right)^{3/8}
\left(r^3\lambda^6\tilde{g}^{-1}\right)^{1/8},
\label{kclqg:eq}
\end{equation}
and
\begin{equation}
k_r = \left(27\ca{C}/44\right)^{3/8}
\left(r^3\lambda^6\tilde{g}^{-1}\right)^{1/8},
\label{kr2lqg:eq}
\end{equation}
where $\tilde{g}$ is the familiar energy generation rate (the rate we
set to unity in the numerical model).  Since the spectral density has
units $\ca{Z} \sim [L^9\tau^{-2}]$, we can also define a spectral
density in unscaled time, $\tilde{\ca{Z}} = \lambda^4\ca{Z}$.  But
$\tilde{g}^{2/3} = \lambda^{4}g^{2/3}$ as well, so that the factors of
$\lambda$ in the unscaled spectra only appear inside the expression
for $k_c$ and $k_r$.  Finally, though, $\ca{Z}$ is a spectrum of
$|\psi|^2/2$, so the APE spectrum is expected to be
\begin{equation}
\ca{APE}(k) = \ca{C}\lambda^2\tilde{g}^{2/3}k^{-11/3}
\left[1-\left(k_c/k\right)^{8/3}\right]^2,
\label{APEspeclqg:eq}
\end{equation}
with $k_c$ given by \eqref{kclqg:eq}.

From \eqref{lqgE:eq}, the spectrum of kinetic energy is the
spectrum of $\ca{E}(k) = k^2\ca{Z}(k)$, so the total KE can be
estimated
\begin{equation}
KE = \int_{k_r}^{\infty}k^2\ca{Z}(k)\, dk
\label{KElqg:eq}
\end{equation}
The rms eddy meridional velocity $V_{\rms}^2 = KE$ (assuming isotropy)
is thus predicted to be
\begin{equation}
V_{\rms}\simeq \left(\frac{33 C^3}{4}\frac{\tilde{g}^3}{r\lambda^2}\right)^{1/8}
\label{Vlqg:eq}
\end{equation}
Following \eqref{diff:eq} the diffusivity of the tracer is then
\begin{equation}
D \simeq \left(\frac{22}{9}\frac{\tilde{g}}{r\lambda^{2}}\right)^{1/2}.
\label{Dlqg:eq}
\end{equation}

Four simulations were performed, differing only in the value of the
linear drag coefficient.  The nondimensional values used were
\[
r^{\star}
= r~\left(gk_f^8\right)^{-1/3}
= r\lambda^2\left(\tilde{g}k_f^8\right)^{-1/3}
=(0.035, 0.23, 1.4, 8.9)\times 10^{-3}.
\]
The spectra of APE, tracer variance and tracer flux, and the
stopping scales, rms velocities and diffusivities, along with
their respective theoretical predictions, are plotted in figure
\ref{lqg:fig}.  At the largest scales, under which the
assumptions used to derive LQG are most accurate, our theory
predicts the spectral shape quite well.  The tracer spectra are
more complicated, but still roughly follow their expected slopes
of $-5/3$ (see discussion of LQG tracer in \S \ref{tracer:sec}).
In the case with the smallest drag, the tracer flux spectrum is
actually \textit{negative} in the transition region $k \lesssim
\lambda$, but we offer no explanation for this behaviour. The
stopping scales and rms velocities are also predicted somewhat
successfully by our theory.  Moreover, the divergence of
simulation from theory is very slight, even when the stopping
scale is on order the deformation scale. The theoretical
prediction for the diffusivity scales like the simulations, but
exceeds the simulated diffusivities by about a factor of 2. Also
shown (the dashed line in the lower right panel of figure
\ref{lqg:fig}) is the prediction for the diffusivity from TDV
with linear drag, which behaves nothing like the simulated
diffusivities.

The physical-space streamfunction and tracer fields for the run with
the lowest drag are shown in the bottom two panels of figure
\ref{phys:fig}.  The streamfunction is quasi-crystalline in structure,
and populated by deformation-scale vortices, as predicted by
\citet{kukharkin95}.  This effect would presumably not be present had
we simulated the strictly defined LQG equation rather than the shallow
water quasi-geostrophic equations.  Despite the predominance of
small-scale ordered structures in the streamfunction, and the
intermittency one would presume to accompany their presence, the
spectrum is well-predicted by the inertial range theory modified for
linear drag.  Moreover, and perhaps more surprisingly, the tracer
field is dominated by smooth plume-like structures whose
characteristic scale is the \textit{energy-containing} scale
associated with the APE, rather than the deformation scale which
strikes the eye in the physical-space image of the streamfunction.

There is a wealth of phenomena in these simulations which require
further research to understand.  The key point of the present
investigation, however, has been to determine the scale responsible
for the mixing of a tracer.

\subsection{Two-dimensional vorticity flow with linear drag and $\beta$}
\label{tdvbr_exp:ssc}

Lastly, we present the results of simulations of \eqref{gen2drb:eq}
with $\alpha=2$.  We have no prediction for the spectral shapes in
this case, since we cannot derive a flux equation similar to
\eqref{epsr:eq} when $\beta$ is present, but we do predict the jet
scale, mixing length and rms velocity at the mixing scale (derived
below). Note that now we have two independent parameters to consider:
$\beta$ and $r$.

In order to see the effects of $\beta$, drag must be set such that
$r\lesssim r_c$, where $r_c$ is defined by \eqref{rcrit2d:eq}.  In the
simulations described here $\beta$ is fixed such that the isotropic
scale \eqref{kbetavm:eq} lies within the computational domain at some
relatively large wavenumber, but below the forcing wavenumber.  The
critical drag $r_c$ is then fixed by $\beta$ and the generation rate
(which is fixed for all simulations described in this paper).  Five
values of drag are then chosen such that the largest value is just at
$r_c$ while the rest are smaller.  In this way we hope to see a
transition from the purely drag-induced diffusivity computed in \S
\ref{tdv_exp:ssc} to the $\beta$ controlled diffusivity, which should
be independent of drag. In particular, we choose
\[
\beta^{\star} \equiv \beta~\left(gk_f^5\right)^{-1/3} = 0.73,  %3450
\]
so that the inviscid $\beta$-scale \eqref{kbetavm:eq} is
(assuming a Kolmogorov constant $\ca{C}=6$ and an effective
generation of $g=g_{\textrm{eff}}\simeq .5$ --- see appendix)
$k_{\beta}/k_{\mathrm{min}}=89.$, or
\[
k_{\beta}/k_f \simeq 0.56.   %89.
\]
The nondimensional critical drag, according to \eqref{rcrit2d:eq}
(using the same values for $g$ and $\ca{C}$ as above), is
\[
r_c^{\star} \simeq 9.0\times 10^{-2}.              %2.1.
\]
The values of drag are then chosen to be
\[
r^{\star} = (0.04, 0.43, 2.1, 4.3, 8.5)\times 10^{-2}
\]
yielding predicted jet-scales from \eqref{kbr:eq}
\[
k_{\beta,r}/k_f \simeq 0.18, 0.33, 0.48, 0.58, 0.69.
\]
% 29., 52., 77., 92., 110.

In figure \ref{kesxy:fig} we try to give a sense of these flows
before showing their statistics.  The upper left panel shows a
contoured snapshot of the physical-space streamfunction for the
case with the lowest drag, and demonstrates that the flow is
dominated by zonal jets, as expected.  The upper right panel
contains the averaged two-dimensional energy spectrum for the
same run.  The surface height is the logarithm of the spectrum and
only the smallest wavenumbers are shown. Energy is concentrated
along the zonal ($k_x=0$) axis, and is excluded from a dumb-bell
region about $k_x=k_y=0$ (an argument for this shape can be found
in \citet{Vallis93a}).  Energy clearly propagates more efficiently
on the zonal axis, where it is uninhibited by the $\beta$-effect.

Slices along the $k_x$ and $k_y$ axes of similar spectra for all five
TDV runs with $\beta$ are shown in the lower left panel of figure
\ref{kesxy:fig}. In order to compare the slopes of these spectra to
our isotropic predictions, remember that integration over angle yields
an integration factor $2\pi k$.  Hence we see that the slopes of the
$k_x$ slice-spectra are about $k^{-8/3}$, which corresponds to an
isotropic spectrum with slope $k^{-5/3}$, while the slopes of the
$k_y$ slice-spectra are about $k^{-6}$, which corresponds to an
isotropic spectrum with slope $k^{-5}$ (see equation
\ref{rossbyspecB:eq}). Both of these slopes are as expected from
arguments made and previous work noted in \S \ref{beta:sec}.
Furthermore, we see that all of the $k_y$ spectra \textit{start to
steepen} at the same scale (the scale at which the $-8/3$ and $-6$
slopes cross), approximately $k/k_f = 0.56 \simeq k_{\beta}/k_f$, as
predicted. (Since $\beta$ is held constant, this scale should be the
same for all simulations.) Interestingly, energy cascades along the
$k_x$ axis past $k\simeq k_{\beta}$, typically peaking at about the
same wavenumber, but with smaller magnitude, than the corresponding
$k_y$ spectra.  This phenomena was also observed by, for example,
\citet{chekhlov96} and \citet{smithwaleffe99}, who point out that
despite the anisotropy present at large scale, an underlying isotropic
cascade still occurs.

A sense of how the two-dimensional spectrum changes as drag is increased
can be gleaned from the bottom right panel of figure \ref{kesxy:fig}, which
shows the inner-most contours of energy values just shy of the maximum
wavenumber along the $k_x$ axis (so just inside the dumb-bell) for the case
with the weakest drag ($r^{\star} = 4\times 10^{-4}$) and the case with the
second to largest drag ($r^{\star} = 4\times 10^{-2}$).  Despite the energy
along the $k_y$ axis starting to steepen at a \textit{fixed} scale for all
five simulations, the dumb-bell increases in size with increasing drag
coefficient.

Let us now consider the statistics of the flow.  The peaks of the energy
spectra along the $k_y$ axis correspond to the jet-scales, which we
expected to be described by \eqref{kbr:eq}.  If the isotropic $\beta$-scale
$k_{\beta}$ is the relevant mixing length, then despite the spectrum of
meridional velocity extending past $k_{\beta}$, we require the meridional
velocity at the mixing scale in order to calculate the diffusivity.  We
estimate the turbulent meridional velocity \textit{near the mixing scale}
as
\begin{equation}
V_{\mathrm{mix}} \simeq \left[\int_{k_{\beta}}^{\infty}\ca{E}(k) \, dk
\right]^{1/2}
\simeq \left(\frac{3}{2}\right)^{1/2}
\left(\frac{\ca{C}^3g^2}{\beta}\right)^{1/5}.
\label{vb:eq}
\end{equation}
Then \eqref{kbetavm:eq} and \eqref{vb:eq} imply a diffusivity
\begin{equation}
D\simeq
\left(\frac{3}{2}\right)^{1/2}
\left(\frac{\ca{C}^{9/2}g^3}{\beta^4}\right)^{1/5},
\label{Dbeta:eq}
\end{equation}
which is independent of $r$.

If, on the other hand, the jet scale $k_{\beta,r}$ or
\eqref{kbr:eq} is the appropriate mixing length, then we expect
$D\sim r^{-1/4}$. Moreover, at values of drag $r\gg r_c$, the
diffusivity should roll over to the slope predicted for TDV
dynamics with linear drag and no $\beta$, namely $D \sim r^{-2}$.

In figure \ref{ldrag_beta:fig} we plot the energy spectra, tracer
variance and tracer flux as functions of total wavenumber, and the
jet-scales (peaks along $k_y$ axis), rms meridional eddy velocities
and integrated tracer fluxes, for the five simulations.  The total
energy spectra show a break at a wavenumber slightly larger than
$k_{\beta}$, where the spectra along the $k_y$ axis begin to steepen.
The shape of these spectra can be best understood by referring back to
the lower left panel of figure \ref{kesxy:fig} --- the contribution to
the spectra in the region with slope $-5$ is almost solely due to
zonal energy.  Strikingly, there is no significant tracer flux at
scales larger than $k_{\beta}$ --- the spectra of the diffusivity (or
tracer flux) vary just slightly as drag is varied, but are essentially
peaked at $k_{\beta}$.  This supports our hypothesis that $k_{\beta}$
is the relevant meridional mixing length.

We also plot some predictions for the jet scales, rms velocities and
diffusivities. The jet-scale is slightly over-predicted by
$k_{\beta,r}$, but for the small drag runs, the slope is close to
$r^{1/4}$, as predicted by \eqref{kbr:eq}.  At values of drag
approaching $r_c$, the slope of jet-scale steepens. The meridional
eddy velocity is somewhat flat at smaller values of $r$, but gently
curves over the entire range.  The fact that the spectrum along the
$k_y=0$ axis increases beyond $k_{\beta}$ in figure \ref{kesxy:fig}
seems to be reflected in the fact that $\overline{v'^2}$ increases
even at the smallest drags examined.

Finally the diffusivity of the tracer, while being quite far off in
magnitude, shows the expected trend: it is approximately constant for
small $r$, then decreases for $r\sim r_c$.  A diffusivity proportional
to $r^{-1/4}$, which would increase by a factor of 4 over the range of
$r$ examined, is not seen, despite the fact that $\overline{v'^2}$
continues to increase slowly with decreasing $r$. In fact, the
diffusivity decreases slightly with decreasing $r$ for the smallest
values examined.  This implies that in the small drag limit, the
inviscid, isotropic $\beta$ scale $k_{\beta}$ of \eqref{kbetavm:eq},
independent of drag, is a reasonable approximation to the mixing
length for the diffusivity estimate.  However, these estimates are
clearly not as accurate in this anisotropic problem, where more
sophisticated closures (e.g., resonant interaction theory) may help in
obtaining quantitatively better fits.


\section{Conclusion}
\label{disc:sec}

This paper has been concerned with the inverse cascade in geostrophic
turbulence, and with the transport of a passive scalar by that
turbulent fluid. We considered conventional two-dimensional (vorticity)
dynamics, surface quasi-geostrophic dynamics, and large-scale
quasi-geostrophic dynamics with a finite deformation radius.

We first showed that, for a general vorticity-streamfunction spectral
field relationship of the form $\xi_k = - k^\alpha \psi$, we may
indeed expect forward and inverse cascades of a generalised enstrophy
and energy. In this context enstrophy refers to the variance of the
advected field, and energy refers to the variance of the associated
velocity. If $\alpha > 0$ ($\alpha < 0$) then generalised energy will
cascade to larger (smaller) scale and generalised enstrophy to smaller
(larger) scale.  [The precise form of this statement is given in
equations \eqref{eq:direc} and \eqref{eq:direcz}.] One notable aspect
of this result is that, if $\alpha < 0$ as for example in large-scale
geostrophic turbulence, the variance of the advected field is
transferred to larger scales, in contrast to the usual
passive-scalar-like behaviour of a cascade to smaller scales.

Given the cascade directions, we can apply classical phenomenology to
obtain predictions of the spectral slope in each inertial range, the
halting scale of the inverse cascade and, at a rather
less-well-founded level, a prediction for a spectrum in the inverse
cascade that includes a frictional decay at small wavenumbers. The
inverse cascades will be modified and made anisotropic by the presence
of a mean gradient of the advected quantity --- the natural
generalization of the familiar beta-effect in TDV. However, because
$\beta$ cannot alter the overall energy level in barotropic flow, the
inverse cascade will only be halted by friction.  [This situation
should be contrasted with the baroclinic case, in which eddy energies
are a highly sensitive function of $\beta$ because of the feedback
between the energy injection rate and the mixing length --- see, \eg
\cite{Held96}.] The anisotropy of the inverse cascade leads to two
distinct length scales: a wavenumber [e.g., \eqref{kbr:eq}]
characterizing the scale of zonal flow, and a smaller scale [e.g.,
\eqref{kbetavm:eq}] that characterizes the mixing length in the
meridional direction.

From these predictions we construct estimates of the diffusion
coefficient of a passive tracer field stirred by the turbulent fluid
and whose variance is maintained by a fixed meridional gradient.  The
expected dependencies for the cases in which only linear drag is
present are
\begin{subequations}
\begin{align}
\textrm{TDV:}
\qquad & D \sim r^{-2}\nonumber\\
\textrm{SQG:}
\qquad & D \sim r^{-1} \ln^{1/2}(k_f/k_c)\nonumber, ~~~
k_c = \left(2\ca{C}/3\right)\left(r^3g^{-1}\right)^{1/3}\\
\textrm{LQG:}
\qquad & D \sim r^{-1/2},\nonumber
\end{align}
\end{subequations}
and, in the small-drag limit when $\beta$ is present,
\begin{subequations}
\begin{align}
\textrm{TDV:}
\qquad & D \sim \beta^{-4/5}\nonumber\\
\textrm{SQG:}
\qquad & D \sim \beta^{-1}\ln^{1/2}(k_f/k_{\beta}), ~~~
k_{\beta} = \beta \ca{C}^{-1/2}g^{-1/3}.\nonumber
\end{align}
\end{subequations}
The last of these follows from using \eqref{kbeta1:eq} for SQG as the
mixing length, but was not derived nor tested here; the LQG case is
special and we do not predict its diffusive behaviour in the presence
of $\beta$, notwithstanding its probable relevance to geophysical
flows. (See \citet{kukharkin96} for some treatment of this case).

In each case in which linear drag halts the cascade, the diffusivity
decreases with increasing drag coefficient.  This is perhaps not
immediately obvious, since, from the perspective of Brownian motion
for example, damping can act to increase the irreversibility of the
flow for a given eddy energy.  In turbulent inverse cascades, however,
decreasing drag allows for increasing eddy scales, thus increasing the
eddy mixing length, and thereby increasing the diffusivity.  On the
other hand, when $\beta$ is significant relative to the drag (\ie when
$r<r_c$ in \eqref{rcrit:eq} or \eqref{rcrit2d:eq}), a diffusivity
which is increasing or constant with drag is reasonable.  In this
less-turbulent, more wave-dominated parameter regime, zonal jets act
as mixing barriers, and the underlying flow is less irreversible.  In
this case, understanding and predicting the transport of tracer likely
requires a different type of mathematical machinery than what we have
employed here.

We have tested many of these predictions numerically. Overall, the
predictions for the halting scale of the inverse cascade, and the
velocity amplitude at the halting scale, are quantitatively
well-satisfied. We saw little evidence of non-universal behaviour over
a wide range of frictional parameters, providing that the inertial
range was well developed and that the halting scale was substantially
smaller than the domain scale, with no pile up of energy at the
halting scale. If any of these are not satisfied the predictions may
fail. (Note also all of our simulations utilized a a single type of
forcing, namely a random stirring localized in wavenumber).  The form
of the predicted spectrum is less well satisfied by the numerical
simulations --- the latter typically falling off much less rapidly
than the former at small, frictionally dominated wavenumbers. However,
because the flow in this low wavenumber regime has little energy, it
does not contribute to the meridional flux of the tracer.

Consistent with these results, the transport of a passive scalar
advected by the turbulent flow is down-gradient diffusive, with a
diffusion coefficient well estimated by turbulence phenomenology.  In
particular, for an isotropic inverse cascade, the diffusion
coefficient is proportional to the product of the mixing length and
the rms velocity amplitude, which are both determined by the halting
scale of the inverse cascade.  In the presence of a beta-effect, the
inverse cascade can develop significant anisotropy; the mixing length
appropriate to the meridional transport of a passive tracer scales
more or less with the isotropic $\beta$-scale given by
\eqref{kbetavm:eq}, whereas the energy containing scale is the
jet-scale \eqref{kbr:eq}.

Our predictions of the diffusion coefficients owe their approximate
success to the presence of a good scale separation between the mean
gradient and the energy containing scales, and a second scale
separation between the energy containing scale and the forcing scale,
enabling a well defined inverse cascade to occur and minimizing the
effects of the forcing.  Both of these conditions are formally
necessary for simple scaling behaviour, but in many cases such
scalings may well remain valid well into more weakly non-linear
regimes \citep[\eg][]{Held96}.  While neither scale separation is
particularly well satisfied by meridional heat transport in the
earth's atmosphere, scale separations between large-scale gradients
and the internal radius of deformation are certainly significant in
the earth's mid-oceans and on the giant gas planets, and the present
diffusivity predictions may be particularly useful in those arenas.

\begin{acknowledgments}
  We acknowledge the contributions of A. Fournier to the tracer
  diffusivity calculations, and thank B. Arbic, S. Danilov, D.
  Gurarie, P. Kushner and G.  Lapeyre for very helpful comments.  S.
  Danilov and D. Gurarie (preprint) have independently proposed and
  explored the scaling \eqref{kbr:eq}.
\end{acknowledgments}

\appendix
        \addtocounter{section}{1}
        \setcounter{equation}{0}

\section*{Appendix A}
\subsection*{Phenomenology of cascade directions}
\label{phenom:sec}

We first define a rescaled wavenumber $k' \equiv k^{\alpha/2}$.
In this space, the energy and enstrophy spectra are such that
\begin{subequations}
\begin{align}
E_G & = \int \ca{E}_G'(k') \,dk' \\
Z_G & =  \int \ca{Z}_G'(k') \,dk'
\end{align}
\end{subequations}
where
\begin{equation}
\ca{Z}_G'(k') = k'^2\ca{E}_G'(k').
\end{equation}
The centroid of the energy distribution in this space is given by
\begin{equation}
\label{eq:centroid2}
k_e' \equiv \frac{\int k' \ca{E}_G'(k')\, dk'}{\int \ca{E}_G'(k') \,dk'}.
\end{equation}
We now suppose that there is some initial distribution of energy
distributed around this wavenumber, and that turbulence
acts irreversibly to broaden this distribution.  That is,
\begin{subequations}
\label{eq:spread3}
\begin{align}
\label{eq:spread3a}
I' & = \int \left( k' - k_e' \right)^2 \ca{E}_G'(k') \, dk' , \\
\frac{d {I'}}{dt} & > 0 .
\end{align}
\end{subequations}
Then, expanding \eqref{eq:spread3a} and  using \eqref{eq:centroid2},
\begin{equation}
  {I'} =  \int (k'^2 - k_e'^2) \ca{E}_G'(k') \, dk' .
\end{equation}
Since $dI'/dt >0$, and both $\int \ca{E}_G'(k') \, dk'$ and  $\int
k'^2 \ca{E}_G'(k') \, dk'$ are constant, we must have that
\begin{equation}
\frac{d k_e'^2}{dt} < 0.
\end{equation}
This implies that
\begin{equation}
\frac{d k^\alpha_e}{dt} < 0,
\label{eq:kdec}
\end{equation}
where $k^\alpha_e$ is a measure of the center of the
energy distribution in conventional wavenumber space, given by
\begin{equation}
k_e \equiv \left[\frac{\int k^{\alpha/2} \ca{E}_G(k)\, dk}
{\int \ca{E}_G(k) \,dk}\right]^{2/\alpha}.
\label{eq:centroid}
\end{equation}
From \eqref{eq:kdec} we immediately have
\begin{equation}
\begin{aligned}
\label{eq:direc}
\frac{d k_e}{dt} & < 0, \qquad \alpha > 0 \\
\frac{d k_e}{dt} & > 0, \qquad \alpha < 0.
\end{aligned}
\end{equation}
It is because the mapping between the two wavenumber spaces is
monotonic that a monotonic direction of energy or enstrophy transfer in
the new wavenumber space corresponds to a monotonic transfer in the
true wavenumber space, although the rates of transfer (and the
direction of transfer if $\alpha < 0$) will differ.

Finally, we note that in the above derivations we may substitute
$\ca{Z}_G(k)$ for $\ca{E}_G(k)$, and $\alpha \rightarrow -\alpha$, in
equations \eqref{eq:centroid} onward. We then readily obtain the
direction of the generalised enstrophy transfer,
\begin{equation}
\begin{aligned}
\label{eq:direcz}
 \frac{d k_z}{dt} & > 0, \qquad \alpha > 0 \\
 \frac{d k_z}{dt} & < 0, \qquad \alpha < 0,
\end{aligned}
\end{equation}
where, analogous to \eqref{eq:centroid},
\begin{equation}
k_z \equiv \left[\frac{\int k^{\alpha/2} \ca{Z}_G(k)\, dk}
{\int \ca{Z}_G(k) \,dk}\right]^{2/\alpha} .
\end{equation}

For TDV and SQG the generalised enstrophy, corresponding to the square
of the advected quantity, cascades to smaller scales and the energy
cascades to larger scales. However, for LQG dynamics the generalised
enstrophy (that is, the variance of the advected quantity) cascades to
\emph{larger} scales, and generalised energy to \emph{smaller}
scales. This occurs in spite of the fact that the equation of motion
is identical to that of a passive tracer, save for the additional
constraint $\Psi = \del^2 \xi$. This serves to emphasize the point
that inferring cascade directions in two-dimensional vorticity
dynamics by analogy to passive tracer dynamics is valid only to the
extent that vorticity acts as a passive tracer. This may well be the
case for small-scale vorticity, but it is not a general rule.

\appendix
        \addtocounter{section}{2}

\section*{Appendix B}
\subsection*{The forcing function}

The forcing is random Markovian in time and normalized at each
timestep such that $g=1$.  Specifically, at each wavenumber $k$ and
timestep $n$, we calculate an initial forcing function
\begin{equation}
\hat{F}_n = cF_{n-1} + (1-c^2)^{1/2}e^{i\theta},
\label{rmf:eq}
\end{equation}
where $\theta$ is a random phase and $c$ is the correlation
coefficient, set to $.99$ in all cases.  Since the generation rate is
$g = -\overline{\psi F}$, in order to make $g=1$ the actual forcing
$F_n$ is calculated as
\begin{equation}
F_n = -\left(\overline{\psi \hat{F_n}}\right)^{-1} \hat{F}_n.
\label{rmf2:eq}
\end{equation}
The effective forcing, however, is less than 1 due to some loss from
the small-scale filter, but this loss is restricted to scales smaller
than the forcing scale.  We calculate this loss and find the effective
generation rate via
\begin{equation}
g_{\mathrm{eff}} = 1 + \langle\overline{\psi \ca{D}_{\psi}}\rangle,
\label{geff:eq}
\end{equation}
where $\ca{D}_{\psi}$ represents the filter dissipation and
$\langle\rangle$ means a time average.  The latter term quickly
reaches a steady value (usually about .45) in all the simulations and
thus the effective rate $g_{\mathrm{eff}}$ is well defined.

\subsection*{The enstrophy filter and timestep}

The equation of motion is stepped forward via a leap-frog step but
with the enstrophy filter applied implicitly,
\begin{equation}
\frac{\xi^{n+1}-\xi^{n-1}}{2\delta_t} = R^n + H(k)\xi^{n+1},
\end{equation}
where $R^n$ represents the right hand side terms at time step $n$ and
$\delta_t$ is the magnitude of the timestep.  Thus
\begin{equation}
\xi^{n+1} =  (2\delta_t R^n + \xi^{n-1})F(k),
\end{equation}
where
\begin{equation}
F(k) = [1 + 2\delta_tH(k)]^{-1}
\end{equation}
is the filter.

While we do not use the following filter, let us, as an example,
consider a hyperviscosity filter of order $s$, \ie
$(-1)^{s/2+1}\nu_s\del^s \xi$, then
\[
H(k) = -\nu_s k^s.
\]
Optimal estimates for the timestep and hyperviscosity are \citep[see,
\eg][]{Maltrud91}
\begin{equation}
\delta_t
\simeq \frac{2\pi}{Nk_{\mathrm{max}}^{2-\alpha}\xi_{\rms}}, ~~~~
\nu_s
\simeq (\xi_{\rms}k_{\mathrm{max}}^{2-\alpha})
k_{\mathrm{max}}^{-s}.
\end{equation}
[The former estimate is used to calculate an adaptive timestep in
the model.] Using these estimates in the hyperviscous filter then
yields
\begin{equation}
F_{\mathrm{hv}}(k) =
\left[1+\frac{4\pi}{N}\left(\frac{k}{k_{\mathrm{max}}}\right)^s\right]^{-1},
\end{equation}
which is independent of the flow.  $N=2k_{\mathrm{max}}$ is the
equivalent horizontal resolution.  Note that $F(0)=1$ and
$F(k_{\mathrm{max}}) = (1+4\pi/N)^{-1}$, so that no dissipation
happens on the largest scale, and the maximum dissipation is a
function only of the resolution.  Furthermore, and sensibly enough, as
$N\rightarrow\infty$, $F(k_{\mathrm{max}})\rightarrow 1$.  Finally,
the power $s$ only affects the curvature of the filter --- the higher
the power, the less affected are the smaller wavenumbers (larger
scales).

Loosely following \citet[][and personal communication]{Lacasce96} we
construct an exponential cutoff filter with the explicit property of
having absolutely no dissipation at wavenumbers smaller than some
preset cutoff wavenumber, $k_{\mathrm{cut}}$.  In particular, suppose
we choose
\begin{equation}
F_{\mathrm{ec}}(k) = \left\{
\begin{array}{ll}
    \exp[-a (k-k_{\mathrm{cut}})^s],     & \mbox{$k>k_{\mathrm{cut}}$}   \\
    1,                                   & \mbox{$k\le k_{\mathrm{cut}}$}.
\end{array}
\right.
\label{fec:eq}
\end{equation}

Let us choose $a$ such that $F(k_{\mathrm{max}}) = (1+4\pi/N)^{-1}$, as
for the optimal hyperviscous filter.  Hence
\begin{equation}
a= \frac{\ln(1+4\pi/N)}{(k_{\mathrm{max}}-k_{\mathrm{cut}})^s}.
\label{a:eq}
\end{equation}
The combination of \eqref{fec:eq} and \eqref{a:eq} make our filter.

Comparisons of the exponential cutoff filter to the hyperviscous
filters are shown for a series of TDV and tracer calculations in
figure \ref{filter_comp:fig}.  Slope of energy spectra with
exponential cutoff filter at small scales is closer to expected $-3$
than for any order of hyperviscous filter.  Also, both tracer variance
and energy are higher at the largest scales with the cutoff filter.

\bibliographystyle{jfm}
%\bibliography{lit}


\begin{thebibliography}{38}
\expandafter\ifx\csname natexlab\endcsname\relax\def\natexlab#1{#1}\fi

\bibitem[Arbic(2000)]{Arbic00}
{\sc Arbic, B.~K.} 2000 Generation of mid-ocean eddies: the local baroclinic
  instability hypothesis. PhD thesis, Massachusetts Institute of Technology.

\bibitem[Bartello \& Holloway(1991)]{Bartello91}
{\sc Bartello, P. \& Holloway, G.} 1991 Passive scalar transport in
  $\beta$-plane turbulence. {\em J. Fluid Mech.\/} {\bf 223}, 521--536.
  
\bibitem[Batchelor(1953)]{Batchelor53} {\sc Batchelor, G.~K.} 1953
  {\em Theory of Homogeneous Turbulence\/}. Cambridge University
  Press.

\bibitem[Borue(1994)]{borue94}
{\sc Borue, V.} 1994 Inverse energy cascade in stationary two-dimensional
  homogeneous turbulence. {\em Phys. Rev. Let.\/} {\bf 72}, 1475--1478.
  
%  \bibitem[Celani {\em et~al.\/}(2001)Celani, Lanotte, Mazzino \&
%  Vergassola]{celani01} {\sc Celani, A., Lanotte, A., Mazzino, A. \&
%  Vergassola, M.} 2001 Fronts in passive scalar turbulence.
%  Preprint.
 
\bibitem[Chekhlov {\em et~al.\/}(1996)Chekhlov, Orszag, Sukoriansky,
  Galperin \& Staroselsky]{chekhlov96} {\sc Chekhlov, A., Orszag,
    S.~A., Sukoriansky, S., Galperin, B. \& Staroselsky, I.} 1996 The
  effect of small-scale forcing on large-scale structures in
  two-dimensional flows. {\em Physica D\/} {\bf 98}, 321--334.
  
\bibitem[Constantin {\em et~al.\/}(1994)Constantin, Majda \&
  Tabak]{constantin94} {\sc Constantin, P., Majda, A.~J. \& Tabak, E.}
  1994 Formation of strong fronts in the 2-d quasigeostrophic thermal
  active scalar. {\em Nonlinearity\/} {\bf 7}, 1495--1533.

\bibitem[Danilov \& Gryanik(2001)]{danilov01b} {\sc Danilov, S. \&
    Gryanik, V.} 2001 Vorticity structure, jets, and spectra in
  quasi-stationary beta-plane turbulence with a strong zonal
  component.  Preprint.
  
\bibitem[Danilov \& Gurarie(2000)]{danilov00} {\sc Danilov, S.~D. \&
    Gurarie, D.} 2000 Quasi-two-dimensional turbulence. {\em Uspekhi
    Fizicheskikh Nauk.\/} {\bf 170}, 921--968.

\bibitem[Danilov \& Gurarie(2001)]{danilov01a}
{\sc Danilov, S. \& Gurarie, D.} 2001 Nonuniversal features of forced
 two-dimensional turbulence in the energy range. {\em Phys. Rev. E\/} {\bf
  63}, 020203:1--4.

\bibitem[Galperin {\em et~al.\/}(2001)Galperin, Sukoriansky \&
  Huang]{galperin01} {\sc Galperin, B., Sukoriansky, S. \& Huang,
    H.~P.} 2001 Universal $n^{-5}$ spectrum of zonal flows on giant
  planets. {\em Phys. Fluids\/} {\bf 13}, 1545--1548.
  
\bibitem[Held \& Larichev(1996)]{Held96} {\sc Held, I.~M. \& Larichev,
    V.~D.} 1996 A scaling theory for horizontally homogeneous,
  baroclinically unstable flow on a beta-plane. {\em J. Atmos.
    Sci.\/} {\bf 53}, 946--952.
  
\bibitem[Held {\em et~al.\/}(1995)Held, Pierrehumbert, Garner \&
  Swanson]{Held95a} {\sc Held, I.~M., Pierrehumbert, R.~T., Garner,
    S.~T. \& Swanson, K.~L.} 1995 Surface quasi-geostrophic dynamics.
  {\em J. Fluid. Mech.\/} {\bf 282}, 1--20.
  
\bibitem[Huang {\em et~al.\/}(2001)Huang, Galperin \&
  Sukoriansky]{huang01} {\sc Huang, H.~P., Galperin, B. \&
    Sukoriansky, S.} 2001 Anisotropic spectra in two-dimensional
  turbulence on the surface of a rotationg sphere. {\em Phys.
    Fluids\/} {\bf 13}, 225--240.

\bibitem[Kolmogorov(1941)]{Kolmogorov41}
{\sc Kolmogorov, A.~N.} 1941 The local structure of turbulence in
  incompressible viscous fluid for very large reynolds numbers. {\em Dokl.
  Acad. Sci. USSR\/} {\bf 30}, 299--303.

\bibitem[Kraichnan \& Montgomery(1980)]{Kraichnan80}
{\sc Kraichnan, R. \& Montgomery, D.} 1980 Two-dimensional turbulence. {\em
  Rep. Prog. Physics\/} {\bf 43}, 547--619.

\bibitem[Kukharkin \& Orszag(1996)]{kukharkin96}
{\sc Kukharkin, N. \& Orszag, S.~A.} 1996 Generation and structure of rossby
  vortices in rotating fluids. {\em Phys. Rev. E\/} {\bf 54}, 4524--4527.

\bibitem[Kukharkin {\em et~al.\/}(1995)Kukharkin, Orszag \&
  Yakhot]{kukharkin95}
{\sc Kukharkin, N., Orszag, S.~A. \& Yakhot, V.} 1995 Quasicrystallization of
  vortices in drift-wave turbulence. {\em Phys. Rev. Lett.\/} {\bf 75},
  2486--2489.

\bibitem[LaCasce(1996)]{Lacasce96}
{\sc LaCasce, J.~H.} 1996 Baroclinic vortices over a sloping bottom. PhD
  thesis, Massachusetts Institute of Technology.

\bibitem[Larichev \& Held(1995)]{Larichev95}
{\sc Larichev, V.~D. \& Held, I.~M.} 1995 Eddy amplitudes and fluxes in a
  homogeneous model of fully developed baroclinic instability. {\em J. Phys.
  Oceanogr.\/} {\bf 25}, 2285--2297.

\bibitem[Larichev \& McWilliams(1991)]{Larichev91}
{\sc Larichev, V.~D. \& McWilliams, J.~C.} 1991 Weakly decaying turbulencs in
  an equivalent-barotropic fluid. {\em Phys. Fluids A\/} {\bf 3}, 938--950.

\bibitem[Leith(1967)]{Leith68a}
{\sc Leith, C.~E.} 1967 Diffusion approximations for two-dimensional
  turbulence. {\em Phys. Fluids\/} {\bf 11}, 671--674.

\bibitem[Lilly(1972)]{Lilly72}
{\sc Lilly, D.~K.} 1972 Numerical simulation studies of two-dimensional
  turbulence. {\em Geophys.\ Astrophys.\ Fluid Dyn.\/} {\bf 3}, 289--319.
  
\bibitem[Lindborg \& Alvelius(2000)]{Lindborg00} {\sc Lindborg, E. \&
    Alvelius, K.} 2000 The kinetic energy spectrum of the
  two-dimensional enstrophy turbulence cascade. {\em Phys. Fluids\/}
  {\bf 12}, 945--947.
  
\bibitem[Maltrud \& Vallis(1991)]{Maltrud91} {\sc Maltrud, M.~E. \&
    Vallis, G.~K.} 1991 Energy spectra and coherent structures in
  forced two-dimensional and beta-plane turbulence. {\em J. Fluid
    Mech.\/} {\bf 228}, 321--342.

\bibitem[Manfroi \& Young(1999)]{manfroi99} {\sc Manfroi, A.~J. \&
    Young, W.~R.} 1999 Slow evolution of zonal jets on the beta plane.
  {\em J. Atmos. Sci.\/} {\bf 56}, 784--800.
  
\bibitem[Oetzel \& Vallis(1997)]{Oetzel97} {\sc Oetzel, K. \& Vallis,
   G.~K} 1997 Strain, vortices, and the enstrophy inertial range in
  two-dimensional turbulence. {\em Phys.  Fluids\/} {\bf 9},
  2991-3004.


\bibitem[Pedlosky(1987)]{Pedlosky87}
{\sc Pedlosky, J.} 1987 {\em Geophysical Fluid Dynamics\/}, 2nd edn. New York:
  Springer.

\bibitem[Pierrehumbert {\em et~al.\/}(1994)Pierrehumbert, Held \&
  Swanson]{Pierrehumbert94}
{\sc Pierrehumbert, R.~T., Held, I.~M. \& Swanson, K.~L.} 1994 Spectra of local
  and nonlocal two-dimensional turbulence. {\em Chaos, Solitons \& Fractals\/}
  {\bf 4}, 1111--1116.

\bibitem[Rhines(1975)]{Rhines75}
{\sc Rhines, P.~B.} 1975 Waves and turbulence on a $\beta$-plane. {\em J.
  Fluid. Mech.\/} {\bf 69}, 417--443.

\bibitem[Rhines(1977)]{Rhines77}
{\sc Rhines, P.~B.} 1977 The dynamics of unsteady currents. In {\em The Sea\/}
  (ed. E.~A. Goldberg, I.~N. McCane, J.~J. O'Briend \& J.~H. Steele), , vol.~6,
  pp. 189--318. J. Wiley and Sons.

\bibitem[Rivera \& Wu(2000)]{Rivera00}
{\sc Rivera, M. \& Wu, X.~L.} 2000 External dissipation in driven
  two-dimensional turbulence. {\em Phys. Rev. Lett.\/} {\bf 5}, 976--979.

\bibitem[Salmon(1980)]{Salmon80}
{\sc Salmon, R.} 1980 Baroclinic instability and geostrophic turbulence. {\em
  Geophys. Astrophys. Fluid Dyn.\/} {\bf 10}, 25--52.

\bibitem[Smith \& Waleffe(1999)]{smithwaleffe99}
{\sc Smith, L.~M. \& Waleffe, F.} 1999 Transfer of energy to two-dimensinoal
  large scales in forced, rotating three-dimensional turbulence. {\em Phys.
  Fluids\/} {\bf 11}, 1608--1622.

\bibitem[Smith \& Yakhot(1993)]{SmithYakhot93}
{\sc Smith, L.~M. \& Yakhot, V.} 1993 Bose condensation and small-scale
  structure generation in a random force driven 2d turbulence. {\em Phys. Rev.
  Let.\/} {\bf 71}, 352--355.

\bibitem[Sukoriansky {\em et~al.\/}(1999)Sukoriansky, Galperin \&
  Chekhlov]{Sukoriansky99}
{\sc Sukoriansky, S., Galperin, B. \& Chekhlov, A.} 1999 Large scale drag
  representation in simulations of two-dimensional turbulence. {\em Phys.
  Fluids\/} {\bf 11}, 3043--3053.

\bibitem[Tennekes \& Lumley(1994)]{Tennekes94}
{\sc Tennekes, H. \& Lumley, J.~L.} 1994 {\em A First Course in Turbulence\/}.
  Cambridge, Massachusetts: MIT Press.

\bibitem[Vallis(1993)]{Vallis93b}
{\sc Vallis, G.~K.} 1993 Problems and phenomenology in two-dimensional
  turbulence. In {\em Nonlinear phenomena in atmospheric and oceanic
  sciences\/} (ed. C.~F. Carnevale \& R.~Pierrehumbert), pp. 1--25. Springer.

\bibitem[Vallis \& Maltrud(1993)]{Vallis93a}
{\sc Vallis, G.~K. \& Maltrud, M.~E.} 1993 Generation of mean flows and jets on
  a beta plane and over topography. {\em J. Phys. Oceanog.\/} {\bf 23},
  1346--1362.

\end{thebibliography}


%**************************************************************************
% Figures

\newcommand{\thecaptionA}{Schematic diagram of the cascades in
\textit{(a)} two-dimensional vorticity dynamics, \textit{(b)} surface
quasi-geostrophic dynamics, and \textit{(c)} large-scale
quasi-geostrophic dynamics. Each panel indicates the cascade
directions, the spectrum of each cascading invariant in its own
inertial range ($\ca{E}_G$ or $\ca{Z}_G$), the spectrum of the
velocity variance ($\ca{E}$), and the dimensional estimate of the
scale at which the inverse cascade is halted by linear drag, $k_r$.
The upward arrows represent the injection of the invariants while the
downward arrows represent their dissipation.}

\newcommand{\thecaptionB}{Equilibrated statistics for five simulations
of TDV dynamics which differ only in the value of the drag
coefficient.  Values of drag used are $r^{\star} = (0.86, 1.5, 2.4,
3.7, 6.0)\times 10^{-2}$ and the effective generation rates are
$g_{\mathrm{eff}} = 0.57, 0.59, 0.61, 0.64, 0.69$.  All statistics are
non-dimensionalized by the timescale \eqref{Tdim:eq} and the length
scale $k_f^{-1}$.  Panels show \textit{(a)} energy spectra as a
function of total wavenumber (solid) and theoretical predictions
(dashed) for each simulation; \textit{(b)} tracer variance spectra as
a function of total wavenumber for each simulation; \textit{(c)} peaks
of energy spectra (stars) and theoretical prediction (solid) vs. drag;
\textit{(d)} rms meridional eddy velocity (stars) and theoretical
prediction (solid) vs. drag; \textit{(e)} tracer flux spectra (the
spectra of $-v'\phi'$) as a function of total wavenumber for each
simulation; \textit{(f)} integrated tracer flux $-\overline{v'\phi'}$
(stars) and theoretical prediction for the diffusivity (solid)
vs. drag.  Theoretical predictions are discussed in \S
\ref{tdv_exp:ssc}.  In this case, the halting scale $k_0$ is compared
to the drag scale $k_r$.}

\newcommand{\thecaptionC}{Snapshots of streamfunction (left) and tracer
(right) fields for TDV (top), SQG (middle) and LQG (bottom)
simulations.  Each field is taken from the equilibrated phase of the
simulation with the lowest drag in each ensemble.  For TDV, $r^{\star}
= 8.6\times 10^{-3}$; for SQG, $r^{\star}= 1.8\times 10^{-3}$; for
LQG, $r^{\star}= 3.5\times 10^{-5}$.}

\newcommand{\thecaptionD}{As for figure \ref{tdv:fig} but for SQG with
varied linear drag.  Values of drag used are $r^{\star} =
(0.18,0.35,0.72,1.4,2.9)\times 10^{-2}$ and the effective generation
rates are $g_{\mathrm{eff}} = 0.35, 0.37, 0.42, 0.54, 0.60$.  The
theoretical predictions for the velocity and diffusivity are slightly
curved due to the relatively large variations in $g_{\mathrm{eff}}$
between simulations.  Theoretical predictions are discussed in \S
\ref{sqg_exp:ssc}.}

\newcommand{\thecaptionE}{Spectra of generalised enstrophy $|\psi|^2$
for LQG with linear drag $r^{\star}= 2.3\times 10^{-4}$ (solid).  Also
shown is theoretical prediction (dashed).  Arrow indicates location of
deformation wavenumber ($\lambda/k_f = 0.44$).}

\newcommand{\thecaptionF}{Same as figure \ref{tdv:fig} but for LQG with
varied linear drag.  Values of the drag coefficient used are
$r^{\star} =(0.035, 0.23, 1.4, 8.9)\times 10^{-3}$ and the effective
generation rates are $g_{\mathrm{eff}} = 0.57, 0.59, 0.62, 0.75$.
Theoretical predictions are discussed in \S \ref{lqg_exp:ssc}.
Missing values in panel \textit{(e)} are negative but small.  Dashed
line in panel \textit{(f)} shows the prediction one would attain for
TDV dynamics with the same flow parameters (see \S
\ref{tdv_exp:ssc}).}

\newcommand{\thecaptionG}{Results from simulations of TDV flow with
fixed $\beta^{\star} = 0.73$ and varied drag.  Values of drag used and
effective generation rates are listed in caption of figure
\ref{ldrag_beta:fig}.  Panels show: \textit{(a)} physical-space
streamfunction contours for case with lowest drag (solid lines are
positive and dashed lines are negative); \textit{(b)} energy spectrum
as a function of $k_x$ and $k_y$ for the simulation with the lowest
drag (surface height is logarithmically scaled energy, and spectrum is
restricted to the $20$ wavenumbers in each direction closest to the
origin); \textit{(c)} zonal ($k_x=0$; dashed) and meridional ($k_y=0$;
solid) slices of energy spectra each simulation; \textit{(d)}
inner-most contours at edge of ``dumb-bell'' region of spectra for
simulations with $r^{\star} = 4.0\times 10^{-4}$ and $4.3\times
10^{-2}$.}

\newcommand{\thecaptionH}{Same as figure \ref{tdv:fig} but TDV flow
with varied drag and fixed $\beta^{\star} = 0.73$.  Values of drag
used are $r^{\star} = (0.040, 0.43, 2.1, 4.3, 8.5)\times 10^{-2}$, and
the critical drag of \eqref{rcrit2d:eq} is $r_c^{\star} = 9.0\times
10^{-2}$.  The effective generation rates are $g_{\mathrm{eff}} =
0.35, 0.37, 0.42, 0.54, 0.60$.  Theoretical predictions are discussed
in \S \ref{tdvbr_exp:ssc}.  Missing values in panel \textit{(e)} are
negative but small.  Dashed lines in bottom panels \textit{(d)} and
\textit{(f)} represent predictions for large-drag, $\beta=0$ limit.}

\newcommand{\thecaptionI}{Comparison of energy spectra for TDV with
various small-scale dissipation functions.  Exponential cutoff filter
(solid) vs. $\del^8$ (dashed), $\del^{12}$ (dash-dot) and $\del^{16}$
(dotted) hyperviscous filters.  Resolution for these runs is $256^2$,
forcing is at $k_f/k_{\mathrm{min}} = 60$ and drag is $r^{\star} =
0.03$.  Also shown are lines representing slopes expected from simple
Kolmogorov scaling.  Exponential cutoff wavenumber in these
simulations is $k_{\mathrm{cut}}/k_{\mathrm{min}} = 70$ and the order
of the filter is $s = 8$ (see \ref{fec:eq} and \ref{a:eq}).}

\clearpage

\begin{figure}
\begin{center}
\vskip 1cm
\includegraphics[width=2.5in]{Figures/schema_tdv.ps}
\put(-190,160){\large \textit{(a)}}
\includegraphics[width=2.5in]{Figures/schema_sqg.ps}
\put(-190,160){\large \textit{(b)}}\\
\includegraphics[width=2.5in]{Figures/schema_lqg.ps}
\put(-190,160){\large \textit{(c)}}
\caption{\thecaptionA}
\label{schema:fig}
\end{center}
\end{figure}

\clearpage

\begin{figure}
\begin{center}
\vskip 1cm
\includegraphics[width=5.5in]{Figures/tdv_ldrag_nc.eps}
\put(-397,523){\large \textit{(a)}}
\put(-185,523){\large \textit{(b)}}
\put(-397,335){\large \textit{(c)}}
\put(-185,335){\large \textit{(d)}}
\put(-397,151){\large \textit{(e)}}
\put(-185,151){\large \textit{(f)}}
\caption{See caption list at end of manuscript.}
\label{tdv:fig}
\end{center}
\end{figure}

\clearpage

\bigskip
\begin{figure}
\begin{center}
\vskip 1cm
\includegraphics[width=2.25in]{Figures/psi_tdv_r1_jet_lr.eps}
\put(-110,170){\textit{Streamfunction} $\psi$}
\put(-190,80){\textit{TDV}}
%\put(-180,150){\large \textit{(a)}}
\hskip 1cm
%\includegraphics[width=2.25in]{Figures/tr_tdv_wmean_r1_jet_lr.eps}
\includegraphics[width=2.25in]{Figures/tr_tdv_r1_jet.eps}
\put(-100,170){\textit{Tracer} $\phi$}
%\put(-180,150){\large \textit{(b)}}\\
\vskip 1cm
\includegraphics[width=2.25in]{Figures/psi_sqg_r023_jet_lr.eps}
\put(-190,80){\textit{SQG}}
%\put(-180,150){\large \textit{(c)}}
\hskip 1cm
%\includegraphics[width=2.25in]{Figures/tr_sqg_r023_wmean_jet_lr.eps}
\includegraphics[width=2.25in]{Figures/tr_sqg_r1_jet.eps}
%\put(-180,150){\large \textit{(d)}}\\
\vskip 1cm
\includegraphics[width=2.25in]{Figures/psi_lqg_jet_lr.eps}
\put(-190,80){\textit{LQG}}
%\put(-180,150){\large \textit{(e)}}
\hskip 1cm
%\includegraphics[width=2.25in]{Figures/tr_lqg_wmean_jet_lr.eps}
\includegraphics[width=2.25in]{Figures/tr_lqg_r1_jet.eps}
%\put(-180,150){\large \textit{(f)}}\\
\vskip 1cm
\caption{\thecaptionC}
\label{phys:fig}
\end{center}
\end{figure}

\clearpage

\begin{figure}
\begin{center}
\vskip 1cm
\includegraphics[width=5.5in]{Figures/sqg_ldrag_nc.eps}
\put(-397,521){\large \textit{(a)}}
\put(-185,521){\large \textit{(b)}}
\put(-397,335){\large \textit{(c)}}
\put(-185,335){\large \textit{(d)}}
\put(-397,151){\large \textit{(e)}}
\put(-185,151){\large \textit{(f)}}
\caption{\thecaptionD}
\label{sqg:fig}
\end{center}
\end{figure}

\clearpage

\begin{figure}
\begin{center}
\vskip 1cm
\includegraphics[width=5.5in]{Figures/lqg_spec_nc.eps}
\caption{\thecaptionE}
\label{lqg_spec:fig}
\end{center}
\end{figure}


\clearpage

\begin{figure}
\begin{center}
\vskip 1cm
\includegraphics[width=5.5in]{Figures/lqg_ldrag_nc.eps}
\put(-397,517){\large \textit{(a)}}
\put(-185,517){\large \textit{(b)}}
\put(-397,330){\large \textit{(c)}}
\put(-185,330){\large \textit{(d)}}
\put(-397,148){\large \textit{(e)}}
\put(-185,148){\large \textit{(f)}}
\caption{\thecaptionF}
\label{lqg:fig}
\end{center}
\end{figure}

\clearpage

\begin{figure}
\begin{center}
\vskip 1cm
\includegraphics[width=5.5in]{Figures/beta_4panel_nc.eps}
\put(-400,390){\large \textit{(a)}}
\put(-185,390){\large \textit{(b)}}
\put(-400,165){\large \textit{(c)}}
\put(-185,165){\large \textit{(d)}}
\vskip 1cm
\caption{\thecaptionG}
\label{kesxy:fig}
\end{center}
\end{figure}

\clearpage

\begin{figure}
\begin{center}
\vskip 1cm
\includegraphics[width=5.5in]{Figures/tdv_beta_ldrag_nc.eps}
\put(-397,517){\large \textit{(a)}}
\put(-185,517){\large \textit{(b)}}
\put(-397,330){\large \textit{(c)}}
\put(-185,330){\large \textit{(d)}}
\put(-397,148){\large \textit{(e)}}
\put(-185,148){\large \textit{(f)}}
\caption{\thecaptionH}
\label{ldrag_beta:fig}
\end{center}
\end{figure}

\clearpage

\begin{figure}
\begin{center}
\vskip 1cm
\includegraphics[width=5.5in]{Figures/filter_comparison_ke_nd.eps}
\caption{\thecaptionI}
\label{filter_comp:fig}
\end{center}
\end{figure}

\clearpage

\noindent{\textbf{\large Figure Captions}}

\noindent \textbf{Figure \ref{schema:fig}}: \thecaptionA

\bigskip
\noindent \textbf{Figure \ref{tdv:fig}}: \thecaptionB

\bigskip
\noindent \textbf{Figure \ref{phys:fig}}: \thecaptionC

\bigskip
\noindent \textbf{Figure \ref{sqg:fig}}: \thecaptionD

\bigskip
\noindent \textbf{Figure \ref{lqg_spec:fig}}: \thecaptionE

\bigskip
\noindent \textbf{Figure \ref{lqg:fig}}: \thecaptionF

\bigskip
\noindent \textbf{Figure \ref{kesxy:fig}}: \thecaptionG

\bigskip
\noindent \textbf{Figure \ref{ldrag_beta:fig}}: \thecaptionH

\bigskip
\noindent \textbf{Figure \ref{filter_comp:fig}}: \thecaptionI

\bigskip


\end{document}
